% Houghton's group H_3 has superpolynomial sofic profile.
% Compile: python build.py (three pdflatex passes).
\documentclass[11pt,a4paper]{article}

\usepackage[utf8]{inputenc}
\usepackage[T1]{fontenc}
\usepackage{lmodern}
\usepackage{amsmath,amssymb,amsthm}
\usepackage[margin=1.1in]{geometry}
\usepackage{booktabs}
\usepackage{enumitem}
\usepackage{microtype}
\setlength{\emergencystretch}{2em}
\usepackage{xcolor}
\usepackage[colorlinks=true,linkcolor=blue!50!black,citecolor=blue!50!black,urlcolor=blue!50!black]{hyperref}
\hypersetup{pdftitle={Houghton's group H3 has superpolynomial sofic profile},
pdfauthor={Seth Douglas and Nidhal Mghirbi},pdfkeywords={sofic profile, hyperlinear profile, Houghton groups, elementary amenable groups, stability, Dehn function}}

\theoremstyle{plain}
\newtheorem{theorem}{Theorem}[section]
\newtheorem{proposition}[theorem]{Proposition}
\newtheorem{corollary}[theorem]{Corollary}
\newtheorem{lemma}[theorem]{Lemma}
\theoremstyle{definition}
\newtheorem{definition}[theorem]{Definition}
\theoremstyle{remark}
\newtheorem{remark}[theorem]{Remark}

\DeclareMathOperator{\Tr}{Tr}
\DeclareMathOperator{\rank}{rank}
\DeclareMathOperator{\diag}{diag}
\DeclareMathOperator{\Sym}{Sym}
\DeclareMathOperator{\FSym}{FSym}

\title{Houghton's group $H_3$ has superpolynomial sofic profile}
\author{Seth Douglas \and Nidhal Mghirbi}
\date{September 2026}

\begin{document}
\maketitle

\begin{abstract}
Amenable groups can be approximated by permutations, but the approximations need not be small. The sofic profile of a finite piece of
a group measures how large a permutation model must be to reproduce its multiplication up to an error $1/r$. We show that Houghton's group
$H_3$, a finitely presented elementary amenable group, has a finite piece whose sofic profile grows at least as fast as $\exp(r^{c})$ with
$c=1/(\log_2130+1)\approx1/8.02$, and at most like $\exp(O(\sqrt r\log r))$. Since $H_3$ is amenable, and hence sofic, this answers the last
question of Cornulier's Problem~3.18: a sofic group can have superpolynomial sofic profile. It also shows that Cornulier's extension theorem,
from which he derived that every elementary amenable group has polynomial sofic profile, is false as stated, and we locate the error: the
finite piece of the kernel that it provides must grow with the error parameter, and we give the corrected, scale-dependent form of the
bound. The same argument bounds unitary approximations, where our unitary models, which spread a defect thinly, give a better upper bound
than our permutation construction. The upper sofic bound holds for every fixed finite chunk of $H_3$.
The proof combines a polynomial bound on the area of far commutators
in Johnson's presentation, obtained from a doubling self-similarity with executable certificates for its finite relations, with an entropy count for
approximately commuting copies of the symmetric group $S_3$. A collection argument also gives the full Dehn-function bound
$\delta_{H_3}(n)=O(n^{\log_2 130+4})$.
\end{abstract}

\tableofcontents

%======================================================================
\section{Introduction}\label{sec:intro}
%======================================================================

Sofic groups are the groups whose finite pieces can be modelled by permutations up to a small error, and hyperlinear groups those
whose finite pieces can be modelled by unitary matrices. Quantitative versions ask, for groups that can be approximated, how large the models must be.
Cornulier introduced the \emph{sofic profile} $\sigma_E(r)$ of a finite piece (a chunk)
$E$ of a group: the least $n$ for which $E$ maps to the symmetric group $\Sym_n$ with multiplication defect $1/r$ and with distinct
elements at normalized Hamming distance at least $1-1/r$~\cite{Cornulier2013}. Amenable groups are sofic. More quantitatively,
Cornulier derived from an extension theorem~\cite[Theorem~3.15]{Cornulier2013} that every elementary amenable
group has polynomial sofic profile~\cite[Example~3.16]{Cornulier2013}, and asked for methods to bound the sofic profile of explicit
groups from below, whether some group has sofic profile that is unbounded and not linear in $r$, and whether such a group can be
sofic~\cite[Problem~3.18]{Cornulier2013}.

Houghton's group $H_3$ is finitely presented and elementary amenable: it is an extension of the locally finite group $\FSym(\mathbb N)$
of finitary permutations by $\mathbb Z^2$~\cite{Houghton1978}. Section~\ref{sec:group} defines it and a chunk $E_H$.

\begin{theorem}\label{thm:sofic}
The chunk $E_H$ of $H_3$ has sofic profile
\[
2^{\,r^{1/(\log130+1)}/940-1/16}\ \le\ \sigma_{E_H}(r)\ \le\ 2^{O(\sqrt r\log r)}\qquad(r>1).
\]
The unweighted bound obtained from certified finite fillings already gives $\sigma_{E_H}(r)\ge2^{\,r^{1/12}/270-1/16}$.
\end{theorem}

Here and below logarithms are to base two, and $\log130\approx7.022$.
For $T\in M_D$ write $\tau_D(T)=D^{-1}\Tr T$,
$\|T\|_2^2=\tau_D(T^*T)$ and $\|T\|_{HS}^2=\Tr(T^*T)$; the dimension subscript on $\tau$ is omitted when clear.
For a density matrix $\rho$, $S(\rho)=-\Tr(\rho\log\rho)$, and
$h_2(p)=-p\log p-(1-p)\log(1-p)$, with $0\log0=0$.

\begin{corollary}\label{cor:intro-cornulier}
Theorem~3.15 of~\cite{Cornulier2013} is false for the extension $1\to\FSym(\mathbb N)\to H_3\to\mathbb Z^2\to1$ and the chunk $E_H$. So $H_3$
is an elementary amenable group whose sofic profile is not polynomial, contrary to~\cite[Example~3.16]{Cornulier2013}.
\end{corollary}

\begin{corollary}[Non-embedding in birational groups]\label{cor:cremona}
For any field $K$ and absolutely irreducible finite-dimensional variety $X$ over $K$, the group $H_3$ does not embed into
$\operatorname{Bir}_K(X)$. The same holds for every group containing $H_3$, in particular $H_m$ for $m\ge3$.
\end{corollary}

\begin{proof}
An embedding preserves the partial multiplication and distinct elements of the chunk $E_H$, and hence its sofic profile.
By the independent Corollary~4.5 of~\cite{Cornulier2013}, every chunk of $\operatorname{Bir}_K(X)$ has polynomial sofic profile;
one may replace $X$ by a nonempty affine open without changing its birational group. This contradicts Theorem~\ref{thm:sofic}.
Finally, $H_3$ embeds into $H_m$ by fixing the additional rays.
\end{proof}

We locate the error in Section~\ref{sec:correction}: the extension theorem asserts a finite piece of the kernel that does not depend on
the error parameter, but the construction it relies on needs kernel elements conjugated along a Følner set whose size grows with the
inverse error. Tracking that dependence gives a correct, scale-dependent extension bound (Proposition~\ref{prop:extension}). Since $H_3$ is
amenable, and hence sofic, Theorem~\ref{thm:sofic} also answers the last question of Cornulier's
Problem~3.18: a sofic group can have sofic profile that is unbounded and not linear in $r$. Slofstra constructed the first finitely
presented group with superpolynomial hyperlinear and sofic profile~\cite{Slofstra2018}, answering the first question; that group is
not amenable, and its soficity and even hyperlinearity were left unresolved there. To our knowledge $H_3$ is the first amenable group shown to have superpolynomial sofic profile.

The lower bound holds for unitary approximations as well. Write $P=\langle a,b,\alpha\mid r_1,\dots,r_5\rangle$ for Johnson's
presentation (Section~\ref{sec:group}) and call unitaries $\phi(a),\phi(b),\phi(\alpha)\in U(D)$ an \emph{$e$-representation} if every
relator satisfies $\|\phi(r_i)-I\|_2\le e$ in the normalized Hilbert--Schmidt norm.

\begin{theorem}\label{thm:unitary}
Let $\beta_H=1/(2\log130+1)\approx1/15.04$, $K'=750\cdot30\cdot3^{\log130}$ and $e>0$. Every $e$-representation of $P$ on $\mathbb C^D$ in which the
three-cycle $\gamma=\alpha\,a\alpha a^{-1}$ satisfies $\|\phi(\gamma)-I\|_2\ge1/2$ has $\log D\ge(K'e)^{-2\beta_H}/16-1/16$.
\end{theorem}

Conversely, spreading the single defect of our permutation models over many cells gives such representations with $\log D=O(e^{-2/3})$,
while our permutation models give only $O(e^{-1})$ (Corollary~\ref{cor:unitary-upper}).

\paragraph{How the proofs go.}
\begin{itemize}[leftmargin=1.2em,itemsep=2pt]
\item \emph{Far commutators have polynomial area} (Section~\ref{sec:area}). The commutator of two conjugates of the basic transposition at
distance $R$ has area at most $30R^{\log130}$ in Johnson's presentation, by a doubling self-similarity of $H_3$ whose finite relations are
certified by machine. Buffered collection upgrades these estimates to a polynomial bound on the full Dehn function
(Theorem~\ref{thm:dehn}), improving the exponential upper bound of~\cite{Lee2012}.
\item \emph{Displaced copies of $S_3$ need dimension} (Section~\ref{sec:profile}). An approximate representation gives, through the far
commutators, many approximately commuting copies of $S_3$ whose three-cycles are displaced from the identity. Rounding each copy to an
exact one and counting the entropy needed to make a state invariant forces dimension exponential in their number.
\item \emph{Transport} (Section~\ref{sec:sofic}). For permutations, relator errors add up in measure over a van Kampen diagram; for
unitaries, in norm. This converts the area bound into the two exponents.
\item \emph{Upper bounds} (Section~\ref{sec:corner}). Permutation models on a torus of clock positions that fail at a single corner, and unitary models that spread
that failure over many cells by small rotations.
\end{itemize}

\paragraph{Related work.} Quantitative metric approximations of groups are surveyed in~\cite{ArzhantsevaCherix2020}; the hyperlinear
profile and its relation to nonlocal games are studied in~\cite{SlofstraVidick2018,Slofstra2018}, and hyperlinear approximations of
amenable groups come from sofic ones~\cite{BCJM2023}. Our method follows Slofstra's in combining stability of approximate
representations of finite groups, here an entropy bound in place of the stability theorems of Gowers--Hatami~\cite{GowersHatami2017}
and De Chiffre--Ozawa--Thom~\cite{DOT2019}, with an isoperimetric input. Houghton's groups have an exponential isoperimetric
inequality~\cite{Lee2012} and are not co-Hopfian~\cite{BCMR2014}; we use the doubling endomorphism behind the latter. The extension
theorem is attributed in~\cite{Cornulier2013} to the construction of Elek and Szabó~\cite{ElekSzabo2006}. A companion paper~\cite{DouglasMghirbiB} applies the
lower bounds to the memory needed for the repeated use of a quantum channel built from $H_3$.

\paragraph{Organization.} Section~\ref{sec:group} defines $H_3$, the chunk and the profiles. Section~\ref{sec:area} proves the area bounds,
Section~\ref{sec:profile} the dimension bounds for approximate representations, and Section~\ref{sec:sofic} the lower bounds of
Theorems~\ref{thm:sofic} and~\ref{thm:unitary}. Section~\ref{sec:corner} gives the upper bounds, Section~\ref{sec:correction} the correction,
and Section~\ref{sec:open} open problems. Appendix~\ref{app:area} gives the fillings.

%======================================================================
\section{Houghton's group, chunks and profiles}\label{sec:group}
%======================================================================

Let $Y=\{(i,j):i\in\{1,2,3\},\,j\ge1\}$ be three rays. Permutations act on the right and
$gh$ means $g$ followed by $h$. Let $a$ send $(1,1)\mapsto(2,1)$, $(1,j)\mapsto(1,j-1)$ for $j\ge2$ and
$(2,j)\mapsto(2,j+1)$, fixing ray~3; let $b$ be defined in the same way with rays 2 and 3 exchanged; let $\alpha$ be
the transposition of $(1,1)$ and $(1,2)$. These generate $H_3$, the group of permutations of $Y$ that are eventually
translations on each ray. They satisfy
\begin{equation}\label{eq:relators}
r_1=\alpha^2,\quad r_2=(\alpha a\alpha a^{-1})^3,\quad r_3=[\alpha,a^2\alpha a^{-2}],\quad
r_4=aba^{-1}b^{-1}\alpha^{-1},\quad r_5=a\alpha a^{-1}b\alpha^{-1}b^{-1},
\end{equation}
with $[x,y]=xyx^{-1}y^{-1}$; each identity can be checked directly on the finitely many points where the word
moves anything. By a theorem of Johnson (see~\cite[Theorem~2.8]{Lee2012}) these relators present $H_3$. The lower
bound below uses only that $r_1,\dots,r_5$ hold in $H_3$, the far-commutator areas of Section~\ref{sec:area}, and
that the three-cycle $\gamma=\alpha\,(a\alpha a^{-1})$ on $(1,1),(1,2),(1,3)$ is nontrivial. The short exact sequence
$1\to\FSym(\mathbb{N})\to H_3\to\mathbb{Z}^2\to1$ shows that $H_3$ is elementary amenable.

We recall Cornulier's definitions~\cite[Definitions~3.2 and~3.4]{Cornulier2013}. A \emph{chunk} of a group $G$ is a finite
subset $E\ni1$ with the partial multiplication $xy=z$ whenever $x,y,z\in E$ and $xy=z$ in $G$. For a map
$f:E\to\Sym_n$ with $f(1)=\mathrm{id}$, let $d_n$ be the normalized Hamming distance on $\Sym_n$. The map is an
\emph{$\varepsilon$-morphism} if $d_n(f(xy),f(x)f(y))\le\varepsilon$ whenever $xy$ is defined in $E$, and it is
\emph{$(1-\varepsilon)$-expansive} if $d_n(f(x),f(y))\ge1-\varepsilon$ for distinct $x,y\in E$. The sofic profile is
\[
\sigma_E(r)=\min\{n:\ \text{there is a }(1-r^{-1})\text{-expansive }r^{-1}\text{-morphism }E\to\Sym_n\},\qquad r>1.
\]
A group has polynomial sofic profile if every chunk has $\sigma_E\preceq$ some polynomial~\cite[Definition~3.10]{Cornulier2013}.

Let $E_H\subset H_3$ be the symmetric chunk consisting of $1$, $a^{\pm1},b^{\pm1},\alpha^{\pm1}$, all prefixes of the relators
$r_1,\dots,r_5$ and of the word $\alpha a\alpha a^{-1}$ for $\gamma$, and the inverses of these elements.

For permutations, relations can be transported in measure: the normalized Hamming distance $d_n$ is bi-invariant, so a
null word that is a product of $A$ conjugates of relators has $d_n(\phi(w),\mathrm{id})\le A\delta$ when every relator has
$d_n(\phi(r_i),\mathrm{id})\le\delta$. For unitaries only $\|\phi(w)-I\|_2\le A\max_i\|\phi(r_i)-I\|_2$ is available in general, since
$\|(e^{i\vartheta}I)^A-I\|_2$ grows like $A\vartheta$.

\section{Polynomial area of far commutators}\label{sec:area}
%======================================================================

Let $P=\langle a,b,\alpha\mid r_1,\dots,r_5\rangle$, put $\sigma_j=a^j\alpha a^{-j}$, and let $A^*(M)$ be the
largest area in $P$ of a commutator $[\sigma_p^{\pm1},\sigma_q^{\pm1}]$ with $0\le p,q\le M$ and $|p-q|\ge2$. These
commutators are trivial in $H_3$ because the transpositions have disjoint supports.

\begin{theorem}[Polynomial far-commutator area]\label{thm:area}
$A^*(M)\le M^{11}$ for every $M\ge2$.
\end{theorem}

\begin{proof}[Outline; the full proof is in Appendix~\ref{app:area}]
Let $\psi$ be the endomorphism of the free group on $a,b,\alpha$ with $\psi(a)=a^2$, $\psi(b)=b^2$ and
$\psi(\alpha)=[a^2,b^2]$. On $H_3$ it is induced by the diagonal action on two copies of $Y$, identified with $Y$ by the bijection
$(1,j,e)\mapsto(1,2j-1+e)$ and $(i,j,e)\mapsto(i,2j-e)$ for $i=2,3$ and $e\in\{0,1\}$.
The resulting group endomorphism, not this bijection of sets, is injective and not surjective
(compare~\cite{BCMR2014}). The proof has four steps.
\begin{enumerate}[label=(\arabic*)]
\item \emph{$\psi$ lifts to $P$ with bounded area multiplier.} Each $\psi(r_i)$ is a null word in $P$, of area at
most $0,17,182,864,359$ for $i=4,1,5,3,2$. Hence $\mathrm{Area}_P(\psi(w))\le864\,\mathrm{Area}_P(w)$.
\item \emph{Two cheap letters.} With $t=b^{-1}a$ and $\rho=a^{-1}b^{-1}ab$, the letter $\alpha$ commutes with $t$ at
cost 1 and with $\rho$ at cost at most 5, and $\{t,\rho\}$ is invariant under $\psi$ up to one relation,
$b^{-1}\rho b=t\rho t^{-1}\cdot\rho\cdot t\rho t^{-1}$, of area at most 14. The identities $\psi(t)=t\rho t$,
$t\rho=b^{-1}tb$ and $\psi(\rho)=t^{-1}\rho^b\rho^{b^2}t\,\rho\,\rho^b$ hold in the free group.
\item \emph{A doubling recursion.} The words $X_n=b^{-n}a^n$ satisfy $X_{2n}=\psi(X_n)$ and $X_{2n+1}=b^{-1}X_{2n}bt$
freely, and in $H_3$ act as a reversal of $[1,n]$ followed by translation on the line formed by rays 2 and 3. Writing
$X_n$ as a word $W_n$ in $t,\rho$ costs $R(n)$ cells, with $R(2n)\le864\,R(n)+84\,\#_\rho(W_n)$ and a comparable odd
step. Unrolling gives $R(n)\le1.35\,n^{\log_2864}$.
\item \emph{From $X_n$ to far commutators.} $[\alpha,X_n]$ expands into one conjugate of $[\alpha,t]^{\pm1}$ or
$[\alpha,\rho]^{\pm1}$ per letter of $W_n$, so its area is at most $2R(n)+5|W_n|$; a far commutator
$[\sigma_0,\sigma_m]$ reduces to it at linear extra cost. This gives $A^*(m)\le5.4\,m^{\log_2864}+10\,m^{\log_2323}+18m\le
m^{11}$ for $m\ge5$, and $A^*(2),A^*(3),A^*(4)\le1,21,47$.
\end{enumerate}
\end{proof}

\paragraph{Certificates.} Every finite relation used in the proof is supplied in machine-checkable form: as an
explicit list of cells $(z_k,r_{i_k},s_k)$ such that the relation equals $\prod_kz_kr_{i_k}^{s_k}z_k^{-1}$ in the free
group on $a,b,\alpha$, which a short independent program verifies by free reduction after recomputing the relation
from its definition. The certified relations and their cell counts are: $A(2)=1$, $A(3)\le21$, $A(4)\le47$,
$A(5)\le79$, $A(6)\le157$; $E(0..3)\le1,4,27,76$; $U(1..3)\le1,5,32$; $B(2),B(3)\le9,12$; the relation of step~(2)
with $14$ cells; $[\alpha,\rho]$ with $5$; $\psi(\alpha)=\sigma_1\sigma_2\sigma_0\sigma_1$ with $6$; the images
$\psi(r_1),\psi(r_2),\psi(r_3),\psi(r_5)$ with $17,359,864,182$; and the rewriting of $\psi(\rho)$ as a word in
$t$ and $\rho$, which uses the relation of step~(2) six times, with $84$. The recursion of step~(3) has also been executed mechanically, transporting certificates through
$\psi$ cell by cell: for $n\le9$ the certificates of $X_n=W_n$ have $0,0,14,84,392,995,2493,5298,11878$ cells, each at
most the recursion bound, and the resulting far commutators satisfy $A(m)\le73\,121$ for $m\le11$, far below $m^{11}$.
Appendix~\ref{app:area} gives the fillings in words; the ancillary files contain the generator, the verifier and the
exported certificates. The exported file holds every finite relation listed above and the recursion for $n\le5$, which the
independent verifier checks; the instances $6\le n\le9$ were executed by the generator but are not in the file, and the generator
exports them for the verifier on request. The proof uses only the finite relations and the hand induction of Appendix~\ref{app:area}.

The same fillings give a better exponent when the five kinds of cells are counted separately.

\begin{theorem}[Reweighted far-commutator area]\label{thm:area-weighted}
$A^*(M)\le30\,M^{\log130}$ for every $M\ge2$.
\end{theorem}

The proof is in Appendix~\ref{app:area}. The same area estimates also control arbitrary null words.

\begin{theorem}[Polynomial Dehn function]\label{thm:dehn}
Let $\delta_P(n)$ be the maximum area of a null word of length at most $n$ in Johnson's presentation. There is an absolute constant $C_0$ such that
\[
\delta_P(n)\le C_0n^4\bigl(1+A^*(4n+2)\bigr)
=O(n^{\log130+4})\qquad(n\ge1).
\]
The unweighted estimate of Theorem~\ref{thm:area} alone gives $\delta_P(n)=O(n^{15})$.
\end{theorem}

The proof in Appendix~\ref{app:dehn} collects a buffered null word into adjacent transpositions, then fills it using a uniform Coxeter insertion algorithm.
This full Dehn bound does not improve the profile exponents: the direct far-commutator estimate is sharper for the words used there.
Reweighting the far-commutator fillings cannot bring their exponent below $\log\rho\approx7.018$, where
$\rho\approx129.64$ is the spectral radius of their substitution matrix; $\log130\approx7.022$ is the certified choice. The area is at least quadratic.

\begin{proposition}[Quadratic lower bound]\label{prop:area-lower}
For odd $k\ge9$, $\mathrm{Area}(C_k)\ge(k-7)(k-5)/8$, where $C_k=[\alpha,a^k\alpha a^{-k}]$. Hence $A^*(R)\ge R^2/16$ for every integer $R\ge23$.
\end{proposition}

\begin{proof}
In a permutation model in which every relator moves at most a fraction $\delta$ of the points, a null word of area $A$ moves at most a
fraction $A\delta$ (transport in measure, Section~\ref{sec:group}). In the corner models of Proposition~\ref{prop:corner}, $\delta\le M^{-2}$
while $C_{2M-1}$ moves a fraction at least $(1-s^{-2})(M-4)(M-3)/(2M^2)$. Hence $\mathrm{Area}(C_{2M-1})\ge(1-s^{-2})(M-4)(M-3)/2$ for every
$s$, which gives the claim with $k=2M-1$. For integer $R\ge23$ apply it to the largest odd $k\le R$.
\end{proof}

Thus the present uniform-area transport and sector estimate cannot give more than the exponent $1/3$ in $r$ for permutation approximations,
and $2/5$ in $1/e$ for unitary approximations with relator defect $e$ (Section~\ref{sec:open}). This is a limit of that argument, not of every
possible use of far commutators.
%======================================================================

\section{Stability and counting}\label{sec:profile}
%======================================================================

The full $S_3^N$ sector bound below feeds the permutation-profile estimate; the phase-robust version later in the section
feeds the unitary estimate and the companion's weighted extraction. The entropy arguments use monotonicity of the Holevo quantity
under partial trace, a consequence of strong subadditivity~\cite{LiebRuskai1973}.

\subsection{An entropic sector bound}

\begin{lemma}[Projection products]\label{lem:projections}
Let $P_1,\dots,P_m$ be orthogonal projections on a Hilbert space and $\|x\|=1$. Then
$1-\mathrm{Re}\langle x,P_m\cdots P_1x\rangle\le2\sum_j\|(I-P_j)x\|^2$.
\end{lemma}

\begin{proof}
Put $x_0=x$, $x_j=P_jx_{j-1}$ and $r_j=x_{j-1}-x_j\in\mathrm{ran}(I-P_j)$. Then $\sum_j\|r_j\|^2=1-\|x_m\|^2=:b$, and
$a=1-\mathrm{Re}\langle x,x_m\rangle=\mathrm{Re}\sum_j\langle(I-P_j)x,r_j\rangle\le\sqrt{Eb}$ with $E=\sum_j\|(I-P_j)x\|^2$. Since
$0\le\|x-x_m\|^2=2a-b$, we get $a\le\sqrt{2Ea}$, that is $a\le2E$.
\end{proof}

\begin{lemma}[Entropic sector bound]\label{lem:sector}
Let $u_1,\dots,u_N:S_3\to U(D)$ be unitary representations, not assumed to commute, such that for a three-cycle $c$
\[
\|u_i(c)-I\|_2\ge\Delta\quad(1\le i\le N),\qquad\|u_i(g)u_j(h)-u_j(h)u_i(g)\|_2\le\eta\quad(i\ne j,\ g,h\in S_3).
\]
Put $t=\sqrt{2N}\,\eta$. If $t\le\Delta$ and $t^2\le5/6$, then
\[
\log D\ \ge\ N\Big[\frac{(\Delta-t)^2}3-h_2(t^2)-t^2\log5\Big].
\]
\end{lemma}

\begin{proof}
Let $G=S_3^N$ and $f(g)=u_1(g_1)\cdots u_N(g_N)$. For $h$ in factor $i$, $\|f(gh)-f(g)u_i(h)\|_2=\|S^*XS-X\|_2$ with $X=u_i(h)$ and
$S=u_{i+1}(g_{i+1})\cdots u_N(g_N)$, and $\|S^*XS-X\|_2^2=2-2\,\mathrm{Re}\,\tau(X^*S^*XS)$. The averages
$T_j(Y)=\frac16\sum_au_j(a)^*Yu_j(a)$ are orthogonal projections on $M_D$ with the normalized Hilbert--Schmidt inner product, the
average of $S^*XS$ over $g$ is $T_N\cdots T_{i+1}(X)$, and $\|(I-T_j)X\|_2^2=\frac1{12}\sum_a\|[X,u_j(a)]\|_2^2\le\eta^2/2$.
Lemma~\ref{lem:projections} therefore gives
\[
\mathbb E_g\|f(gh)-f(g)u_i(h)\|_2^2\le2(N-i)\eta^2\le t^2 .
\]
Let $R$ be the right regular representation of $G$ and $W:\mathbb C^D\to\ell^2(G)\otimes\mathbb C^D$ the isometry
$Wv=|G|^{-1/2}\sum_g|g\rangle\otimes f(g)v$. Then $\|R(h)W-Wu_i(h)\|_{HS}\le\sqrt D\,t$ for $h$ in factor $i$. Put $P=WW^*$ and
$\rho=P/D$, a state of entropy $\log D$. For the three-cycle $c_i$ of factor $i$ and $Z_i=2I-R(c_i)-R(c_i)^*$,
\[
\mathrm{Tr}(\rho Z_i)=D^{-1}\|(R(c_i)-I)W\|_{HS}^2\ge(\Delta-t)^2 .
\]
Let $\mathcal E_i(X)=\frac16\sum_hR(h_i)XR(h_i)^*$, with $h_i$ the element $h$ placed in factor $i$; these conditional expectations
commute. Since $(I-P)W=0$, $\mathrm{Tr}[(I-P)\mathcal E_i(\rho)]\le t^2$. The state $\mathcal E_i(\rho)$ is an average of six states of rank
$D$, one of them $\rho$, so after pinching by $P$ and $I-P$ its two blocks have rank at most $D$ and $5D$. Pinching does not decrease
entropy and $x\mapsto h_2(x)+x\log5$ increases on $[0,5/6]$, so $S(\mathcal E_i(\rho))\le\log D+h_2(t^2)+t^2\log5$. For a conditional
expectation $S(\mathcal E(\sigma))-S(\sigma)=D(\sigma\|\mathcal E(\sigma))$; since the $\mathcal E_j$ commute, monotonicity of relative entropy
under $\mathcal E_{i-1}\cdots\mathcal E_1$ bounds the entropy increment of $\mathcal E_i$ applied after $\mathcal E_{i-1}\cdots\mathcal E_1$ by
$D(\rho\|\mathcal E_i\rho)$. Hence $\bar\rho=\mathcal E_N\cdots\mathcal E_1(\rho)$ satisfies $S(\bar\rho)\le\log D+N[h_2(t^2)+t^2\log5]$.

The state $\bar\rho$ commutes with $R(G)$, so $\bar\rho=\bigoplus_\lambda p_\lambda(I_{d_\lambda}/d_\lambda\otimes\sigma_\lambda)$ over the irreducible
representations $\lambda$ of $S_3^N$, and $S(\bar\rho)\ge\sum_\lambda p_\lambda\log d_\lambda=\sum_\lambda p_\lambda k_\lambda$, where $k_\lambda$ is the
number of factors of $\lambda$ carrying the two-dimensional representation of $S_3$. The operator $Z_i$ is central in the group
algebra, so $\mathrm{Tr}(\bar\rho Z_i)=\mathrm{Tr}(\rho Z_i)$, and it acts as $3I$ on the sectors whose factor $i$ is two-dimensional and as
$0$ on the others. Hence $N(\Delta-t)^2\le\sum_i\mathrm{Tr}(\bar\rho Z_i)=3\sum_\lambda p_\lambda k_\lambda\le3S(\bar\rho)$.
\end{proof}

\begin{lemma}[Exact $S_3$ rounding]\label{lem:s3round}
If $A,B\in U(D)$ satisfy $\|A^2-I\|_2,\|B^2-I\|_2,\|(AB)^3-I\|_2\le e$, there is a unitary representation $u$ of
$S_3=\langle s,t\mid s^2,t^2,(st)^3\rangle$ on $\mathbb C^D$ with $\|u(s)-A\|_2\le e$, $\|u(t)-B\|_2\le8e$ and $\|u(st)-AB\|_2\le9e$, and
each normal form $1,s,t,st,ts,sts$ evaluated in $u$ is within $10e$ of the same word in $A,B$.
\end{lemma}

\begin{proof}
Round the spectra of $A$ and $B$ to the nearer of $\pm1$. For $|z|=1$ this moves $z$ by at most $|z^2-1|$, so the resulting
Hermitian unitaries $S,T$ satisfy $\|S-A\|_2,\|T-B\|_2\le e$. Then $C=ST$ satisfies $SCS=C^*$ and $\|C^3-I\|_2\le7e$. Round the spectrum
of $C$ to the nearest cube root of unity, choosing $1$ at ties and at $-1$. This rule commutes with complex conjugation and moves
$z$ by at most $|z^3-1|$, so $C_0$ satisfies $C_0^3=I$, $SC_0S=C_0^*$ and $\|C_0-C\|_2\le7e$. Then $u(s)=S$ and $u(t)=SC_0$ are
involutions with $u(s)u(t)=C_0$, and the bounds follow from the triangle inequality.
\end{proof}

\subsection{A phase-robust form}

The count of Lemma~\ref{lem:sector} can be run on products of three-cycles alone. It then needs no relation between the
transpositions of different copies, so a coherent phase in their commutators does no harm. Write $\kappa(u,v)=\|uv-vu\|_2$ and
$G(t)=h_2(\min\{\frac23t,\frac23\})+\min\{\frac23t,\frac23\}$.

\begin{theorem}[Phase-robust sector bound]\label{thm:phase-robust}
Let $c_i,x_i\in U(D)$, $1\le i\le N$, satisfy $c_i^3=x_i^2=I$ and $x_ic_ix_i=c_i^{-1}$, with no condition on the commutators of
$x_i$ and $x_j$. Put $\Delta_i=\|c_i-I\|_2$ and $f(g)=c_1^{g_1}\cdots c_N^{g_N}$ for $g\in\mathbb Z_3^N$, let $g^{(i)}$ be $g$ with its
$i$-th coordinate negated, and, for uniform $g$,
\[
\hat a_i^2=\mathbb E_g\|f(g+e_i)-f(g)c_i\|_2^2,\qquad \hat b_i^2=\mathbb E_g\|x_if(g^{(i)})-f(g)x_i\|_2^2 .
\]
Then $\hat a_i^2\le\frac43\sum_{j>i}\kappa(c_i,c_j)^2$, $\hat b_i^2\le\frac43\sum_{j\ne i}\kappa(x_i,c_j)^2$, and
\[
\log D\ \ge\ \sum_{i=1}^N\Big[\frac{(\Delta_i^2-\hat a_i^2)_+}3-\frac{\hat b_i^2}{2\ln2}-G(\hat a_i^2)\Big].
\]
In particular $\Delta_i\ge49/100$ and $\hat a_i,\hat b_i\le1/100$ for all $i$ give $\log D\ge N/16$.
\end{theorem}

\begin{proof}
\emph{The ordered-product errors.} For $h$ unitary, $\|(I-E_j)h\|_2^2=\frac13\kappa(h,c_j)^2$ for the average $E_j$ of conjugation by the
powers of $c_j$. Moving $c_i$ through the suffix $c_{i+1}^{g_{i+1}}\cdots c_N^{g_N}$ and averaging over its exponents gives
$\hat a_i^2=2-2\,\mathrm{Re}\langle c_i,E_N\cdots E_{i+1}c_i\rangle$, and Lemma~\ref{lem:projections} gives the first bound. For the second,
$T_j(Z)=\frac13\sum_kc_j^{-k}Zc_j^{s_jk}$ with $s_i=-1$ and $s_j=1$ otherwise are orthogonal projections, $T_i(x_i)=x_i$ by the local
relation, $\|(I-T_j)x_i\|_2^2=\frac13\kappa(x_i,c_j)^2$ for $j\ne i$, and $\hat b_i^2=2-2\,\mathrm{Re}\langle x_i,T_N\cdots T_1x_i\rangle$.

\emph{Embedding.} Let $V\xi=3^{-N/2}\sum_gf(g)\xi\otimes|g\rangle$, let $C_i$ shift the $i$-th coordinate of $\ell^2(\mathbb Z_3^N)$ by $-1$, let
$F_i$ negate it, and put $X_i=x_i\otimes F_i$. The $C_i$ commute, $X_iC_iX_i=C_i^{-1}$, $X_iC_jX_i=C_j$ for $j\ne i$, and
$D^{-1/2}\|C_iV-Vc_i\|_{HS}=\hat a_i$, $D^{-1/2}\|X_iV-Vx_i\|_{HS}=\hat b_i$.

\emph{Entropy.} Let $\rho=VV^*/D$ and let $\mathcal T_i$ average conjugation by the powers of $C_i$. The twirl moves weight at most
$\min\{\frac23\hat a_i^2,\frac23\}$ outside the range of $VV^*$, into a block of rank at most $2D$, so $S(\mathcal T_i\rho)\le\log D+G(\hat a_i^2)$; as in
Lemma~\ref{lem:sector} the twirls commute, and the increment of each twirl, the Holevo quantity of the conjugates of the current
state, can only decrease under the earlier twirls by data processing for the Holevo quantity, so $\bar\rho=\mathcal T_1\cdots\mathcal T_N\rho$ has
$S(\bar\rho)\le\log D+\sum_iG(\hat a_i^2)$.

\emph{Sectors.} Let $p$ be the joint distribution of the eigenvalue labels $Z_i\in\{1,\omega,\omega^2\}$ of the $C_i$, the same for $\rho$ and
$\bar\rho$. Then $S(\bar\rho)\ge H(p)\ge\sum_iH(Z_i\mid Z_{-i})$. Let $Q_a$ and $P_b$ be the spectral projections of $c_i$ and $C_i$. The
numbers $r_{ab}=D^{-1}\|P_bVQ_a\|_{HS}^2$ form a joint distribution whose marginals are the spectral distribution of $c_i$ in the
normalized trace and that of $C_i$ in $\rho$. Distinct cube roots of unity are at squared distance $3$, so $\hat a_i^2=3\Pr(a\ne b)$ and
$\Delta_i^2=3\Pr(a\ne1)$, whence $s_i:=\Pr(Z_i\ne1)\ge(\Delta_i^2-\hat a_i^2)_+/3$. Next, the normalized vectorizations of $V$ and of
$X_iVx_i^*$ are at squared distance $\hat b_i^2$, so their real overlap is $1-\hat b_i^2/2$. Measuring the joint spectral projections of
the $C_j$ gives $p$ from the first and its flip $p^{(i)}$ in coordinate $i$ from the second, because $X_i$ exchanges $\omega$ and $\omega^2$
in coordinate $i$ and fixes the other labels; by the Cauchy--Schwarz inequality outcome by outcome,
$\sum_z\sqrt{p(z)p(z^{(i)})}\ge1-\hat b_i^2/2$. For $0\le t\le1$ one has $1-h_2(t)\le(1-2\sqrt{t(1-t)})/\ln2$: with $y=2t-1$ the difference of
$1-\sqrt{1-y^2}$ and $\ln2\,(1-h_2(t))$ vanishes with its derivative at $y=0$ and has second derivative
$(1-y^2)^{-3/2}-(1-y^2)^{-1}\ge0$. So for fixed other labels with masses $u,w$ on the two nontrivial values of $Z_i$, the conditional
entropy receives at least $u+w-(u+w-2\sqrt{uw})/\ln2$, and summing gives $H(Z_i\mid Z_{-i})\ge s_i-\hat b_i^2/(2\ln2)$. The numbers in the
last claim give $\frac{0.49^2-10^{-4}}3-\frac{10^{-4}}{2\ln2}-G(10^{-4})>\frac1{16}$.
\end{proof}

The theorem isolates the row quantities $\hat a_i,\hat b_i$ that control the dimension estimate used below.

\subsection{From an approximate presentation to dimension}

Call $\phi:\{a,b,\alpha\}\to U(D)$ an $e$-representation if $\|\phi(r_i)-I\|_2\le e$ for $i=1,\dots,5$, where $\phi$
is extended to the free group. If $w$ is null in $P$ with area $A$, then $\|\phi(w)-I\|_2\le Ae$, because $w$ is freely a
product of $A$ conjugates of relators and the norm is unitarily invariant.

\begin{proposition}\label{prop:dimension}
Let $\phi$ be an $e$-representation with $\|\phi(\gamma)-I\|_2\ge1/2$, and let $N\ge1$ satisfy $750\sqrt N\,A_Ne\le1$, where
$A_N=\max\{1,A^*(3N)\}$. Then $\log D\ge N/16$.
\end{proposition}

\begin{proof}
Put $A_i=\phi(\sigma_{3(i-1)})$ and $B_i=\phi(\sigma_{3(i-1)+1})$ for $1\le i\le N$; all indices are at most $3N-2$. In $H_3$ each pair
generates a copy of $S_3$ on $\{(1,3i-2),(1,3i-1),(1,3i)\}$; its involution and braid relations are conjugates of $r_1$ and $r_2$, with
defect at most $e$. Lemma~\ref{lem:s3round} gives exact local pairs $x_i=u_i(s)$ and $c_i=u_i(st)$ with $\|x_i-A_i\|_2\le e$ and
$\|c_i-A_iB_i\|_2\le9e$, and $A_iB_i=\phi(a)^{3(i-1)}\phi(\gamma)\phi(a)^{-3(i-1)}$, so $\Delta_i\ge\frac12-9e$. Pairs with different $i$ have
indices differing by at least two, so four and two far-commutator interchanges give $\kappa(c_i,c_j)\le4A_Ne+36e\le40A_Ne$ and
$\kappa(x_i,c_j)\le2A_Ne+20e\le22A_Ne$. By Theorem~\ref{thm:phase-robust} and the concavity of $G$,
$\log D\ge N[(\overline{\Delta^2}-\bar a)/3-\bar b/(2\ln2)-G(\bar a)]$ with $\bar a=N^{-1}\sum_i\hat a_i^2\le\frac{2\cdot40^2}{3\cdot750^2}$,
$\bar b=N^{-1}\sum_i\hat b_i^2\le\frac{4\cdot22^2}{3\cdot750^2}$ and $\overline{\Delta^2}\ge(\frac12-\frac9{750})^2$. Exact rational arithmetic, using
$\ln2>69/100$ and $h_2(q)\le q\log(1/q)+q/\ln2$, gives a margin of at least $141043/1397250000$ above $1/16$.
Explicitly, put $a=2\cdot40^2/(3\cdot750^2)$, $b=4\cdot22^2/(3\cdot750^2)$ and $q=2a/3=64/50625$.
The integer inequality $64^3 2^{29}\ge50625^3$ gives $\log(1/q)\le29/3$, and the resulting rational lower margin is
\[
\frac{(488/1000)^2-a}{3}-\frac{b}{2(69/100)}
-\left(\frac{32}{3}+\frac{100}{69}\right)q-\frac1{16}
=\frac{141043}{1397250000}>0.
\]
\end{proof}

If instead $\|\phi(\gamma)-I\|_2\ge\Delta$ with $0<\Delta\le1$, the same argument with $750$ replaced by a constant depending only on $\Delta$
gives $\log D\ge c_\Delta N$ with $c_\Delta>0$. With $A^*(M)\le CM^p$ the condition reads $N^{p+1/2}e=O(1)$, which gives
$\log D\ge c\,e^{-2/(2p+1)}$: the exponent is $2/(2p+1)$ in $1/e$, or $1/(2p+1)$ in the squared defect $1/e^2$. For $p=\log130$ this is the
exponent $2\beta_H$ of Theorem~\ref{thm:unitary}.

%======================================================================
\section{Lower bounds}\label{sec:sofic}
%======================================================================

\begin{proposition}\label{prop:dimension-perm}
Let $\phi$ assign permutations of a finite set $\Omega$ to $a,b,\alpha$, with $d_n(\phi(r_i),\mathrm{id})\le\delta$ for all $i$ and
$d_n(\phi(\gamma),\mathrm{id})\ge1/2$, and let $N\ge1$ satisfy $10^8N\,A_N\delta\le1$, where $A_N=\max\{1,A^*(3N)\}$. Then
$\log|\Omega|\ge N/16$.
\end{proposition}

\begin{proof}
Use the pairs $A_i=\phi(\sigma_{3(i-1)})$, $B_i=\phi(\sigma_{3(i-1)+1})$, $1\le i\le N$, as permutation matrices, with $\|P-I\|_2^2=2d_n(P,\mathrm{id})$. Their
$S_3$ relators have $d_n\le\delta$, hence $\|\cdot\|_2$-defect at most $e=\sqrt{2\delta}$, and Lemma~\ref{lem:s3round} gives exact
representations $u_i$ within $10e$ of the normal forms. Commuting two normal forms from different pairs takes at most nine
far-commutator interchanges, so by transport in measure their commutator has $d_n\le9A_N\delta$ and $\|\cdot\|_2\le3\sqrt{2A_N\delta}$;
rounding adds at most $40e$, so $\eta\le43\sqrt{2A_N\delta}$. The element $u_i(st)$ lies within $9e$ of a conjugate of $\phi(\gamma)$,
whose $\|\cdot\|_2$-displacement is at least $1$, so $\Delta\ge1-9\sqrt{2\delta}\ge\frac12$. Then $t^2=2N\eta^2\le7396NA_N\delta<10^{-4}$, and
Lemma~\ref{lem:sector} gives $\log|\Omega|\ge N[(49/100)^2/3-h_2(10^{-4})-10^{-4}\log5]\ge N/16$.
\end{proof}

With $A^*(M)\le CM^p$ the condition reads $N^{p+1}\delta=O(1)$, which gives the sofic exponent $1/(p+1)$.

\begin{proof}[Proof of Theorem~\ref{thm:sofic}]
Let $f:E_H\to\Sym_n$ be a $(1-r^{-1})$-expansive $r^{-1}$-morphism with $r\ge10$, and regard permutations as unitary matrices.
For permutation matrices, $\|P-P'\|_2^2=2\,d_n(P,P')$, and $d_n$ is bi-invariant. Put $\phi(s)=f(s)$ for $s\in\{a,b,\alpha\}$
and extend $\phi$ to the free group. Since $s\cdot s^{-1}=1$ in the chunk, $d_n(f(s)^{-1},f(s^{-1}))\le r^{-1}$. For a relator
$s_1\cdots s_l$ with prefixes $p_k$, the chunk contains $p_ks_{k+1}=p_{k+1}$ and $p_l=1$, so $l-1$ merges and at most $l$
inverse replacements give $d_n(\phi(r_i),\mathrm{id})\le(2l-1)/r\le23/r$, that is $\|\phi(r_i)-I\|_2\le\sqrt{46/r}$. In the
same way $d_n(\phi(\gamma),f(\gamma))\le4/r$, while expansiveness gives $d_n(f(\gamma),\mathrm{id})\ge1-1/r$ because $\gamma\ne1$;
so $d_n(\phi(\gamma),\mathrm{id})\ge1-5/r\ge1/2$. Apply Proposition~\ref{prop:dimension-perm} with $\delta=23/r$ and
$A^*(3N)\le30(3N)^{\log130}$ (Theorem~\ref{thm:area-weighted}): the condition $10^8N\cdot30(3N)^{\log130}\delta\le1$ holds for
$N=\lfloor(r/(23\cdot10^8\cdot30\cdot3^{\log130}))^{1/(\log130+1)}\rfloor$. Therefore, whenever $N\ge1$ (which forces $r\ge10$),
\[
\log\sigma_{E_H}(r)\ \ge\ \frac N{16}\ \ge\ c_Sr^{1/(\log130+1)}-\frac1{16},
\]
with $c_S=(23\cdot10^8\cdot30\cdot3^{\log130})^{-1/(\log130+1)}/16>1/940$,
and for smaller $r$ the right side is negative. For unitary approximations Proposition~\ref{prop:dimension} gives the
corresponding bound: by Theorem~\ref{thm:area-weighted}, with $K'=750\cdot30\cdot3^{\log130}$ and $N=\lfloor(K'e)^{-2\beta_H}\rfloor$, every
$e$-representation of $P$ on $\mathbb C^D$ with $\|\phi(\gamma)-I\|_2\ge1/2$ has $\log D\ge(K'e)^{-2\beta_H}/16-1/16$, where
$\beta_H=1/(2\log130+1)$; this is Theorem~\ref{thm:unitary}. The upper bound
$\log\sigma_{E_H}(r)=O(\sqrt r\log r)$ is Proposition~\ref{prop:sofic-upper}.
\end{proof}

\section{Upper bounds}\label{sec:corner}

\subsection{Permutation models}

The upper bounds use permutation models that follow the quotient $\mathbb Z^2$ on a torus of clock positions and fail at a single
corner. Fix integers $M\ge2$ and $s\ge2$, put $\ell=2M+2$ and $L=3\ell$, and let $\Omega_{M,s}=(\mathbb Z/M\mathbb Z)^2\times[s]^L$.
For $2\le r\le L$ let $c_r=(1\ r\ r{-}1\ \cdots\ 2)\in\Sym(L)$ be the reverse cycle and $\tau=(1\ 2)$; products of permutations are read
from left to right, and a permutation $\pi$ of the coordinates acts on colorings by $(x\cdot\pi)(t)=x(t\pi^{-1})$. For $0\le i,j<M$ put
$A_{i,j}=c_{\ell-i-j}$ and, for $j<M-1$, $B_{i,j}=c_{2\ell-j}$. On the top row put $B_{i,M-1}=D_i$, where
\begin{equation}\label{eq:corner-rec}
D_0=c_{2\ell-M+1},\qquad D_{i+1}=c_{\ell-i-(M-1)}^{-1}\,\tau\,D_i\,c_{\ell-i}\qquad(0\le i<M-1).
\end{equation}
Define permutations of $\Omega_{M,s}$, with clock coordinates modulo $M$, by
\[
(i,j,x)\cdot a=(i+1,j,x\cdot A_{i,j}),\qquad(i,j,x)\cdot b=(i,j+1,x\cdot B_{i,j}),\qquad(i,j,x)\cdot\alpha=(i,j,x\cdot\tau).
\]

\begin{proposition}[Corner models]\label{prop:corner}
In these permutations $r_1,r_2,r_3,r_5$ hold exactly and $r_4$ fails only at the clock $(M-1,M-1)$, where it acts as the product
$R_M=(1\ \,M{+}4)\prod_{t=5}^{M+3}(t\ \,t{+}M)$ of $M$ disjoint transpositions of the colours; so every relator moves at most
a fraction $M^{-2}$ of $\Omega_{M,s}$, and $\gamma$ moves a fraction $1-s^{-2}$. For $M\ge5$ the far commutator
$C_{2M-1}=[\alpha,a^{2M-1}\alpha a^{1-2M}]$ moves a fraction at least $(1-s^{-2})(M-4)(M-3)/(2M^2)$. A word of length at most $t$ that
represents a nonidentity element of $H_3$ fixes at most a fraction $s^{-1}+4t/M$ of $\Omega_{M,s}$ when $M>2t$.
\end{proposition}

\begin{proof}
\emph{Cycle identities.} For $r\ge4$, $c_r\tau c_r^{-1}=(2\ 3)$ and $c_r(2\ 3)c_r^{-1}=(3\ 4)$, and for $3\le r<t$
\begin{equation}\label{eq:corner-id}
c_rc_tc_{r-1}^{-1}c_t^{-1}=\tau ;
\end{equation}
both sides exchange $1$ and $2$ and fix every other point. By induction on $i$, the recursion~\eqref{eq:corner-rec} gives $D_i(t)=t-1$
for $2\le t\le M+3-i$: this holds for the reverse cycle $D_0$, and a point $t\le M+2-i$ passes the four factors of $D_{i+1}$ as
$t\mapsto t+1\mapsto t+1\mapsto t\mapsto t-1$. Since $M+3-i\ge4$, $D_i^{-1}(1)=2$ and $D_i^{-1}(2)=3$, so $D_i\tau D_i^{-1}=(2\ 3)$.

\emph{Relators.} At every clock the colour permutation of $a\alpha a^{-1}$ is $(2\ 3)$, that of $a^2\alpha a^{-2}$ is $(3\ 4)$, and that of
$b\alpha b^{-1}$ is $(2\ 3)$, on the top row by the property of $D_i$. Hence $r_1,r_2,r_3,r_5$ act trivially everywhere, and $\gamma$ acts at
every clock as a three-cycle of coordinates, which fixes a fraction $s^{-2}$ of the colorings. At the clock $(i,j)$ the colour
permutation of $r_4$ is $A_{i,j}B_{i+1,j}A_{i,j+1}^{-1}B_{i,j}^{-1}\tau$. For $j<M-1$ it is trivial by~\eqref{eq:corner-id} with
$r=\ell-i-j$ and $t=2\ell-j$, also for $i=M-1$ because $B_{i,j}$ does not depend on $i$ off the top row. For $j=M-1$ and $i<M-1$ it is
trivial by~\eqref{eq:corner-rec}. Only the clock $(M-1,M-1)$ remains. There the colour permutation is $c_4D_0c_{M+3}^{-1}D_{M-1}^{-1}\tau$. With
$E_r=c_r^{-1}\tau=(2\ 3\ \cdots\ r)$ the recursion unrolls to $D_{M-1}=E_5E_6\cdots E_{M+3}\,D_0\,c_{2M+2}c_{2M+1}\cdots c_{M+4}$, and
$D_0c_r^{-1}D_0^{-1}=E_{r+1}$ for $r\le3M+4$, so the corner permutation is $c_4\,H\,G^{-1}\tau$ with $H=E_{M+4}\cdots E_{2M+3}$ and
$G=E_5\cdots E_{M+3}$. Following the points, $H$ sends $t\mapsto t+M$ for $2\le t\le M+3$ and $t\mapsto2M+5-t$ for
$M+4\le t\le2M+3$, and $G$ sends $t\mapsto t+M-1$ for $2\le t\le4$ and $t\mapsto M+5-t$ for $5\le t\le M+3$, both fixing all other
points. Hence $HG^{-1}=(2\ 3\ 4\ \,M{+}4)\prod_{t=5}^{M+3}(t\ \,t{+}M)$, and $c_4HG^{-1}\tau=R_M$.

\emph{Far commutators.} Let $k=2M-1$, fix a clock $(i,j)$ with $i+j\ge M+3$ (so $i\ge4$), and put $F=\ell-j$. The colour permutation
of $a^k$ read from this clock is $P=\prod_{t=0}^{k-1}c_{F-((i+t)\bmod M)}$, and that of $\sigma_k$ is $P\tau P^{-1}$, which moves
$P^{-1}(1)$ and $P^{-1}(2)$. With $d_r=c_r^{-1}=(1\ 2\ \cdots\ r)$,
\[
x\cdot(d_rd_{r+1}\cdots d_t)=\begin{cases}x+t-r+1,&x<r,\\ t-x+1,&r\le x\le t,\\ x,&x>t.\end{cases}
\]
The factors of $P^{-1}$ come in three blocks, with cycle lengths $F-i+2,\dots,F$, then $F-M+1,\dots,F$, then $F-M+1,\dots,F-i$. The
first block sends $(1,2)$ to $(i,i+1)$; since $i\ge F-M+1$, the second sends $(i,i+1)$ to $(F-i+1,F-i)$; the third fixes $F-i+1$ and
sends $F-i$ to $1$. So $\sigma_k$ acts as the transposition $(1\ \ 2M+3-i-j)$, whose second point is at least $5$, and $C_k$ acts as
a three-cycle. There are $(M-4)(M-3)/2$ such clocks.

\emph{Short words.} Let $\mathcal S_{i,j}$ consist of the first $\ell-i-j$, $\ell+i$ and $\ell+j$ points of rays $1,2,3$, ordered ray by
ray; it has $L$ points. When the clock step does not wrap, the generator $a$ of $H_3$ maps $\mathcal S_{i,j}$ onto $\mathcal S_{i+1,j}$ and
induces $c_{\ell-i-j}$ on the ordered coordinates; likewise $b$ induces $c_{2\ell-j}$ and $\alpha$ induces $\tau$. So on the clocks at
distance at least $t$ from the lines where the clock wraps, a fraction at least $1-4t/M$ when $M>2t$, a word of length at most $t$
acts through the element of $H_3$ it represents. If that element has nonzero translation it fixes no state. If it is nontrivial
with zero translation, it moves only points among the first $t+2$ of each ray, which lie in $\mathcal S_{i,j}$, so it acts as a
nonidentity permutation of the coordinates and fixes at most a fraction $s^{-1}$ of the colorings. Its fixed-point fraction is
therefore at most $s^{-1}+4t/M$.
\end{proof}

\begin{proposition}[Upper bound for every fixed chunk]\label{prop:sofic-upper}
For every finite chunk $E\subset H_3$, $\log\sigma_E(r)=O_E(\sqrt r\log r)$. In particular this holds for $E_H$.
\end{proposition}

\begin{proof}
Fix representative words for the elements of $E$, with the empty word representing $1$, and let $t$ bound the lengths of their pairwise
differences. The finitely many chunk relations have some maximum area $B_E$ in Johnson's presentation. Take $s=M$ and the diagonal action
on two independent copies of $\Omega_{M,M}$. For $M>4t$, the short-word estimate in Proposition~\ref{prop:corner} bounds the fixed-point fraction
of every nonidentity difference by $((1+8t)/M)^2$ in the product: the fractions square. Each relator has defect at most $2/M^2$ in the product,
so transport in measure bounds every chunk-relation defect by $2B_E/M^2$.
Choosing $M$ at least a sufficiently large $E$-dependent multiple of $\sqrt r$ gives an $r^{-1}$-morphism that is $(1-r^{-1})$-expansive.
Its degree is $|\Omega_{M,M}|^2=M^{4+2(6M+6)}$, whose logarithm is $O(M\log M)$, proving the claim.
\end{proof}

\subsection{Unitary models}\label{sec:unitary-models}

In a permutation model a relator that fails costs the fraction of points it moves. In a unitary model it costs a sum of squared
angles, so a single defect can be spread thin. Doing this for the corner models gives unitary models whose relator defect has squared norm of order $M^{-3}$, not $M^{-2}$, at the same size. Throughout, products are read from left to right, the left factor acting first,
as in Section~\ref{sec:corner}, and $e_1=(1,0)$, $e_2=(0,1)$.

\emph{One-particle models.} No generator moves the last $2M+2$ of the $6M+6$ colour coordinates. Use the reduced colour coordinates $1,\dots,L$ with $L=4M+4$, and read each link permutation $A_c$, $B_c$
and $\tau$ of Section~\ref{sec:corner} as a permutation matrix on $\mathbb C^L$. A one-particle model replaces some of these link matrices
by unitaries. For a word $w$ and a clock $c$ the holonomy $W_c(w)\in U(L)$ is the product of the link matrices along the path of $w$ from
$c$, with inverse matrices for inverse letters. For a two-element set $S=\{x,y\}$ let $Q_S$ be the projector onto
$(e_x-e_y)/\sqrt2$; then $\exp(i\pi Q_S)$ is the transposition $(x\ y)$, and $\exp(i\vartheta Q_S)$ is a phase rotation in the plane of $S$.

\emph{Fermionic quantization.} For $W\in U(L)$ let $\Lambda(W)$ act on the exterior algebra $\Lambda(\mathbb C^L)$, of dimension $2^L$,
by $v_1\wedge\cdots\wedge v_m\mapsto Wv_1\wedge\cdots\wedge Wv_m$. It is multiplicative on all of $U(L)$, and
$\mathrm{Tr}\,\Lambda(W)=\det(I+W)$, since in degree $m$ its eigenvalues are the products of $m$ eigenvalues of $W$. For $k\ge1$ put
$\Lambda_k=\Lambda^{\otimes k}$, so that $2^{-kL}\,\mathrm{Tr}\,\Lambda_k(W)=\det\big((I+W)/2\big)^k$. The quantized model acts on $\mathcal H=\mathbb C^{M^2}\otimes\Lambda(\mathbb C^L)^{\otimes k}$, with normalized trace $\tau_{\mathcal H}$, distinct from the transposition $\tau=(1\ 2)$. A letter moves the clock $c$ to
$c+\pi(s)$ and applies $\Lambda_k$ of its link matrix at $c$. These right-action operators compose in reverse order; taking their adjoints defines the usual free-group representation $\rho$ on $a,b,\alpha$. A word whose translation is nonzero
modulo $M$ has $\tau_{\mathcal H}(\rho(w))=0$, and a word that closes has
\[
\|\rho(w)-I\|_2^2=\mathbb E_c\Big[2-2\,\mathrm{Re}\det\big((I+W_c(w))/2\big)^k\Big],
\]
with $\mathbb E_c$ the average over clocks. For a three-cycle $\sigma$ of coordinates $\det((I+\sigma)/2)=\frac14$.

\begin{lemma}[Rotating the links]\label{lem:rotations}
In a corner model, replace every link $A_c$ by $U_cA_c$ and every link $B_c$ by $V_cB_c$, where $U_c$ commutes with $(2\ 3)$ and
$(3\ 4)$, and $V_c$ commutes with $(2\ 3)$. Then at every clock the holonomies of $\alpha$, $a\alpha a^{-1}$, $a^2\alpha a^{-2}$ and
$b\alpha b^{-1}$ are $(1\ 2)$, $(2\ 3)$, $(3\ 4)$ and $(2\ 3)$, so $r_1,r_2,r_3,r_5$ have identity holonomy everywhere and $\gamma$ acts
as a three-cycle.
\end{lemma}

\begin{proof}
By the proof of Proposition~\ref{prop:corner}, $A_c\tau A_c^{-1}=(2\ 3)$, $A_c(2\ 3)A_c^{-1}=(3\ 4)$ and $B_c\tau B_c^{-1}=(2\ 3)$ at every
clock. So the holonomy of $a\alpha a^{-1}$ at $c$ is $U_c(2\ 3)U_c^{-1}=(2\ 3)$, that of $a^2\alpha a^{-2}$ is
$U_cA_c\,U_{c+e_1}(2\ 3)U_{c+e_1}^{-1}A_c^{-1}U_c^{-1}=U_c(3\ 4)U_c^{-1}=(3\ 4)$, and that of $b\alpha b^{-1}$ is $V_c(2\ 3)V_c^{-1}=(2\ 3)$.
The relators $r_1,r_2,r_3,r_5$ are products of these conjugates, as in the proof of Proposition~\ref{prop:corner}.
\end{proof}

A plaquette $q$ with base $v_0$ has the vertices $v_0$, $v_1=v_0+e_1$, $v_2=v_0+e_2$, $v_0+e_1+e_2$; its bottom and top edges are the
$a$-links at $v_0$ and $v_2$, and its left and right edges the $b$-links at $v_0$ and $v_1$. Its $r_4$ holonomy is
$W_q=A_{v_0}B_{v_1}A_{v_2}^{-1}B_{v_0}^{-1}\tau$. A \emph{region} consists of a set $R$ of plaquettes, a source $q_0\in R$, a sign
$s_R=\pm1$, and a two-element label $\mathrm{lab}(v)\subset\{1,\dots,L\}$ at every vertex of every plaquette of $R$. An edge is
\emph{interior} if it lies on two plaquettes of $R$.

\begin{lemma}[Spreading a defect]\label{lem:spreading}
Take a permutation model satisfying the local identities
$A_c\tau A_c^{-1}=(2\ 3)$, $A_c(2\ 3)A_c^{-1}=(3\ 4)$ and $B_c\tau B_c^{-1}=(2\ 3)$ of Lemma~\ref{lem:rotations} at every clock.
Suppose $W_q$ is the identity except at finitely many sources, where it is the
product of the transpositions given by the base labels of the regions with that source. Let finitely many regions satisfy:
\begin{enumerate}[label=(R\arabic*),itemsep=0pt]
\item every edge of a non-source plaquette, from $v$ to $w$, carries $\mathrm{lab}(v)$ onto $\mathrm{lab}(w)$ under its link permutation;
\item every base label avoids $\{1,2\}$, and the label at the initial vertex of every interior edge avoids $\{2,3,4\}$;
\item each region is connected through its interior edges, and the only interior edge of its source is the bottom edge;
\item labels of different regions are disjoint wherever both are defined; on an interior link carrying a region's rotation, its label is disjoint from any additional modification of that link.
\end{enumerate}
Here the link permutations belong to the base model, including any preliminary exact moves. Potentials vanish on boundary links;
no disjointness from an additional modification is required for a region on a link where that region carries no rotation.
Let $F_{R,q}$ be real numbers with $\sum_{q\in R}F_{R,q}=s_R\pi$ for each region. Then there are potentials $\varphi_R$ on the interior
edges of $R$ such that the model with $U_c=\prod_R\exp(i\varphi_R(\text{$a$-edge at }c)\,Q_{\mathrm{lab}_R(c)})$ and
$V_c=\prod_R\exp(i\varphi_R(\text{$b$-edge at }c)\,Q_{\mathrm{lab}_R(c)})$ still has $r_1,r_2,r_3,r_5$ exact, and its $r_4$ holonomy at
every plaquette $q$ is $\prod_{R\ni q}\exp(iF_{R,q}Q_{\mathrm{lab}_R(v_0)})$, with $v_0$ the base of $q$.
\end{lemma}

\begin{proof}
Put $\varphi_R=0$ off the interior edges and $(\mathrm{curl}\,\varphi)(q)=\varphi(\text{bottom})+\varphi(\text{right})-\varphi(\text{top})
-\varphi(\text{left})$. An interior edge enters the curls of its two plaquettes with opposite signs, so the curls over $R$ sum to zero.
Choose a spanning tree of the plaquettes of $R$ joined through interior edges, rooted at the source, and fix the potential of each tree edge,
from the leaves inward, so that $\mathrm{curl}\,\varphi_R=F_R-s_R\pi\,\delta_{q_0}$ at the plaquette below it; since both sides sum to zero,
the equation then also holds at the root. By (R2) and (R4) all rotations on a link commute with the transpositions of
Lemma~\ref{lem:rotations}, so $r_1,r_2,r_3,r_5$ stay exact.

For the $r_4$ holonomy, the unmodified identity $A_{v_0}B_{v_1}A_{v_2}^{-1}B_{v_0}^{-1}=W_q\tau$ gives
\[
W_q'=U_{v_0}\big(A_{v_0}V_{v_1}A_{v_0}^{-1}\big)\,W_q\,\tau\big(B_{v_0}U_{v_2}B_{v_0}^{-1}\big)^{-1}V_{v_0}^{-1}\tau .
\]
Conjugation by a link permutation moves the plane of a rotation, and at a non-source plaquette (R1) moves the planes $\mathrm{lab}(v_1)$ and
$\mathrm{lab}(v_2)$ to $\mathrm{lab}(v_0)$. At a source only the bottom edge carries a potential. So all rotations in $W_q'$ lie in the
planes $\mathrm{lab}_R(v_0)$; they avoid $\{1,2\}$ by (R2), hence commute with $\tau$, and planes of different regions are disjoint by (R4).
For region $R$ the angles add up to $(\mathrm{curl}\,\varphi_R)(q)$, and at its source $W_q$ contributes the rotation by $s_R\pi$ in the
plane of the source label. The total is $F_{R,q}$.
\end{proof}

\begin{proposition}[Unitary models]\label{prop:unitary-models}
For every $M\ge48$ the model above with $k=2$ and suitable rotations of the links gives unitaries $\rho(a),\rho(b),\rho(\alpha)$ on a space of
dimension $M^2\,2^{8M+8}$ in which $r_1,r_2,r_3,r_5$ act as the identity, $\|\rho(r_4)-I\|_2^2\le14.26\,M^{-3}$ and $\|\rho(\gamma)-I\|_2^2=15/8$.
\end{proposition}

\begin{proof}
\emph{Two exact moves.} Let $c^*=(M-1,M-1)$, where the $r_4$ holonomy is $R_M$ (Proposition~\ref{prop:corner}). Replace $A_{c^*}$ by
$\sigma_1A_{c^*}$ and $B_{(0,M-2)}$ by $\sigma_2B_{(0,M-2)}$, with $\sigma_1=(1\ \,M{+}4)$ and $\sigma_2=(5\ \,M{+}5)$. These satisfy the
hypotheses of Lemma~\ref{lem:rotations}, so $r_1,r_2,r_3,r_5$ stay exact. The modified links lie on three plaquettes. At $c^*$ the
holonomy becomes $\sigma_1R_M=\prod_{t=5}^{M+3}(t\ \,t{+}M)$. At $(0,M-2)$, where it was the identity, it becomes
$\tau\sigma_2\tau=\sigma_2$. At $(M-1,M-2)$ it becomes $(A\sigma_2A^{-1})\,\tau(B\sigma_1B^{-1})\tau$ with $A=A_{(M-1,M-2)}=c_5$ and
$B=B_{(M-1,M-2)}=c_{3M+6}$. Here $A\sigma_2A^{-1}=(1\ \,M{+}5)$ and $B\sigma_1B^{-1}=(2\ \,M{+}5)$, so $\tau(B\sigma_1B^{-1})\tau=(1\ \,M{+}5)$ and
the holonomy stays the identity. Use these modified permutations as the base model in Lemma~\ref{lem:spreading}.

\emph{The regions.} For $5\le t\le M+3$ let $R_t$ have source $c^*$ and sign $(-1)^{t-5}$, and consist of $c^*$, the plaquettes $(M-1,j)$
with $1\le j\le M-2$, and the plaquettes $(i,j)$ with $0\le i\le M-3$, $0\le j\le M-2$ and $i+j\le t+M-8$; label the vertex $(i,j)$ by
$\{t+M-1-i-j,\ t+2M-1-j\}$ for $0\le i\le M-2$ and $(M-1,j)$ by $\{t+M-1-j,\ t+2M-1-j\}$. Let $T$ have source $(0,M-2)$ and sign
$(-1)^{M-1}$, and consist of $(0,M-2)$ and the plaquettes $(i,j)$ with $i,j\ge0$ and $i+j\le M-3$, with labels $\{M+3-i-j,\ 2M+3-j\}$.
The source labels are the defect pairs $\{t,t+M\}$ at $c^*$ and $\{5,M+5\}$ at $(0,M-2)$.

These regions satisfy (R1)--(R4). Off the top row $B_{i,j}=c_{4M+4-j}$ lowers every label coordinate by one, which is the change of the
labels from $j$ to $j+1$. In the triangles $A_{i,j}=c_{2M+2-i-j}$ lowers the first label coordinate, which is at most $2M+2-i-j$ since
$t\le M+3$, and fixes the second, which exceeds $2M+2-i-j$ since $t\ge5$; this is the change from $i$ to $i+1$, and the same holds in $T$.
On the column $i=M-1$, $A_{M-1,j}=c_{M+3-j}$ fixes both label coordinates, which are at least $t+M-1-j\ge M+4-j$, and the labels at
$(M-1,j)$ and $(0,j)$ agree; this is the only crossing of a seam, since the rows stay below $M-1$ and $c^*$ is attached only through its
bottom edge. Every base label and every label at the initial vertex of an interior edge has both coordinates at least $5$. Two labels of
bulk regions at one vertex differ in $t$ by at most $M-2$, so they are disjoint; a bulk label meets a label of $T$ only if
$t\in\{4,\ 4-M-i,\ M+4+i\}$, which is impossible; and $\sigma_1$, $\sigma_2$ are disjoint from the labels $\{t,t+M\}$ and $\{t+1,t+M+1\}$
on the links where those regions carry rotations. At $(0,M-2)$ the label of $T$ equals the support of $\sigma_2$, but that $b$-link is a boundary link of $T$ and carries no $T$-rotation.
Each region is connected through interior edges: the triangle is a staircase, the column is glued to it through
the $b$-edges at $(0,j)$, and each source through its bottom edge. Finally $R_t\subset R_{t+1}$, and every region has at least
$m_M=(M-1)(M-2)/2+1$ plaquettes.

\emph{The defect.} Take $F_{R,q}=s_R\pi/|R|$ on each region and $k=2$. By Lemma~\ref{lem:spreading} the $r_4$ holonomy at $q$ is a product
of phase rotations by angles $F_{R,q}$ in disjoint planes, so $\det((I+W'_q)/2)=\prod_R(1+e^{iF_{R,q}})/2=e^{i\Theta_q}C_q$ with
$\Theta_q=\frac12\sum_RF_{R,q}$ and $C_q=\prod_R\cos(F_{R,q}/2)\in[0,1]$. Using $1-C^k\le k(1-C)$, $1-\prod_Rc_R\le\sum_R(1-c_R)$ and
$1-\cos x\le x^2/2$,
\[
2-2\,\mathrm{Re}\big(e^{ik\Theta_q}C_q^k\big)=2(1-C_q^k)+2C_q^k(1-\cos k\Theta_q)\le\frac k4\sum_RF_{R,q}^2+k^2\Theta_q^2 .
\]
Summed over the plaquettes, $\sum_q\sum_RF_{R,q}^2=\sum_R\pi^2/|R|\le M\pi^2/m_M$. The bulk regions containing $q$ are those with $t$
at least some $t_q$; their signs alternate and their magnitudes $\pi/|R_t|$ decrease, so they sum to at most $\pi/m_M$ in absolute value,
and $T$ adds at most $\pi/m_M$. Hence $|\Theta_q|\le\pi/m_M$, and averaging over the $M^2$ clocks,
\[
M^3e_4^2:=M^3\|\rho(r_4)-I\|_2^2\le\frac{\pi^2M^2}{2m_M}+\frac{4\pi^2M^3}{m_M^2}
\le\frac{\pi^2M^2}{(M-1)(M-2)}+\frac{16\pi^2M^3}{(M-1)^2(M-2)^2},
\]
which decreases in $M$ and is below $14.26$ at $M=48$. So $e_4^2\le14.26\,M^{-3}$ for $M\ge48$.

\emph{The three-cycle.} By Lemma~\ref{lem:rotations}, $\gamma$ acts at every clock as a three-cycle $\sigma$ of coordinates, so
$\tau(\rho(\gamma))=\det((I+\sigma)/2)^2=4^{-2}$ and $\|\rho(\gamma)-I\|_2^2=2-2\cdot4^{-2}=15/8$.
\end{proof}

\begin{corollary}[Unitary upper bound]\label{cor:unitary-upper}
Let $0<e\le e_0=14.26^{1/2}\cdot48^{-3/2}$ and $M=\lceil(14.26/e^2)^{1/3}\rceil$. Johnson's presentation has an $e$-representation on $\mathbb C^D$ with
$\|\phi(\gamma)-I\|_2\ge1/2$ and $\log D\le8M+8+2\log M$. So such representations exist with $\log D=O(e^{-2/3})$.
\end{corollary}

\begin{proof}
The choice of $M$ gives $M\ge48$ and $14.26\,M^{-3}\le e^2$; apply Proposition~\ref{prop:unitary-models}.
\end{proof}

Our permutation construction gives the weaker bound in this norm. Read as permutation matrices, the corner models with $s=2$ have relator
defects at most
$\sqrt2/M$ in $\|\cdot\|_2$, since a permutation moving a fraction $f$ of the points has $\|P-I\|_2^2=2f$, and $\|\phi(\gamma)-I\|_2^2=3/2$, at
$\log D=6M+6+2\log M$; this gives only $\log D=O(e^{-1})$. Spreading the defect lowers the exponent of $1/e$ in our upper bounds from $1$ to
$\frac23$; whether optimal permutation models are weaker than optimal unitary ones is open (Section~\ref{sec:open}). The lower
bound of Theorem~\ref{thm:unitary} has exponent $2\beta_H\approx0.133$.

%======================================================================
\section{The correction}\label{sec:correction}
%======================================================================

\begin{corollary}\label{cor:cornulier}
Theorem~3.15 of~\cite{Cornulier2013} is false for $1\to\FSym(\mathbb{N})\to H_3\to\mathbb{Z}^2\to1$ and the chunk $E_H$.
Consequently the conclusion of~\cite[Example~3.16]{Cornulier2013}, that every elementary amenable group has polynomial
sofic profile, fails for $H_3$.
\end{corollary}

\begin{proof}
Theorem~3.15 of~\cite{Cornulier2013} asserts a symmetric chunk $E'$ of the kernel and a finite symmetric $S\subset\mathbb{Z}^2$
with $\sigma_{E_H}(r)\le\sigma_{E'}(r)\,\alpha_{\mathbb{Z}^2,S}(r)$ for all $r>1$, where $\alpha_{Q,S}(r)$ is the least size of a
finite $F\subset Q$ with $|SF\setminus F|<|F|/r$. The kernel $\FSym(\mathbb{N})$ is locally finite, so $E'$ lies in a finite
subgroup and $\sigma_{E'}$ is bounded~\cite[Fact~3.5]{Cornulier2013}; boxes show $\alpha_{\mathbb{Z}^2,S}(r)=O(r^2)$. The resulting
bound $\sigma_{E_H}(r)=O(r^2)$ contradicts Theorem~\ref{thm:sofic}. Since $H_3$ is elementary amenable, Example~3.16 fails. The
contradiction does not need Theorem~\ref{thm:area-weighted}: with the certified bound of Theorem~\ref{thm:area} alone, the proof of
Theorem~\ref{thm:sofic} gives $\sigma_{E_H}(r)\ge2^{\,r^{1/12}/270-1/16}$.
\end{proof}

\paragraph{Where the argument breaks.} Theorem~3.15 is stated in~\cite{Cornulier2013} as what the explicit construction
of Elek and Szabó~\cite{ElekSzabo2006} yields ``without any change''. That construction, given a finite $F\subset G$ and
$\varepsilon>0$, chooses a Følner set $A\subset Q$ with $|Ag\setminus A|\le\varepsilon|A|$ for $g\in F$, a section
$\sigma:Q\to G$, and then an approximation of the kernel on the finite set
\[
H_\varepsilon=N\cap\big(\sigma(A)\,F\,\sigma(A)^{-1}\big),
\]
acting on $B\times A$ with $B$ the kernel's approximation space~\cite[proof of Theorem~1(3)]{ElekSzabo2006}. The kernel
chunk is therefore $H_\varepsilon$, and it depends on $\varepsilon=1/r$ through $A$. For $H_3$ with the section
$(m,n)\mapsto a^mb^n$ and $A=[0,M)^2$, it contains $a^mb^n\alpha b^{-n}a^{-m}$, the transposition of $(1,m+n+1)$ and
$(1,m+n+2)$, for all $0\le m,n<M$: at least $2M-1$ distinct transpositions, which generate the symmetric group of $2M$
points, with $M$ of order $r$. No chunk independent of $r$ contains them. What the construction does give is the bound of
Proposition~\ref{prop:extension} below, in which the kernel chunk grows with $r$. The upper bound by the isoperimetric
profile~\cite[Proposition~3.14]{Cornulier2013} is not affected, and $H_3$, having exponential growth, has isoperimetric profile at least
exponential. The proof of~\cite[Proposition~3.17]{Cornulier2013} for Baumslag--Solitar groups applies Theorem~3.15 in the same way, so it is
incomplete as well; we do not know whether its conclusion holds.

\paragraph{A scale-dependent extension bound.} The construction of Elek and Szabó, with the accuracy tracked explicitly, gives the following
correct form of the extension bound. Let $1\to N\to G\xrightarrow{\pi}Q\to1$ be an extension, $s:Q\to G$ a section with $s(e)=1$, and put
$c(a,g)=s(a)\,g\,s(a\pi(g))^{-1}\in N$ for $a\in Q$ and $g\in G$.

\begin{proposition}\label{prop:extension}
Let $E\subset G$ be a symmetric chunk, $r>1$, and $A\subset Q$ a finite nonempty set with $|A\pi(g)\setminus A|\le|A|/(4r)$ for every $g\in E$. Then
$K_{A,E}=\{c(a,g):\ g\in E,\ a\in A,\ a\pi(g)\in A\}$ is a symmetric chunk of $N$, and
\[
\sigma_E(r)\ \le\ |A|\,\sigma_{K_{A,E}}(2r).
\]
\end{proposition}

\begin{proof}
The set $K_{A,E}$ contains $c(a,1)=1$, and $c(a,g)^{-1}=c(a\pi(g),g^{-1})$, so it is symmetric. The cocycle identity
$c(a,g)\,c(a\pi(g),h)=c(a,gh)$ holds in $G$. Let $f_K:K_{A,E}\to\Sym_m$ be a $(1-\frac1{2r})$-expansive $\frac1{2r}$-morphism with $m=\sigma_{K_{A,E}}(2r)$, and
on $A\times[m]$ let $f(g)$ send $(a,x)$ to $(a\pi(g),x\cdot f_K(c(a,g)))$ whenever $a\pi(g)\in A$, completed arbitrarily to a permutation; $f(1)$ is the
identity. Call a fiber $a$ good for a relation $gh=k$ in $E$ if $a$, $a\pi(g)$ and $a\pi(k)$ lie in $A$. Then $c(a,g)c(a\pi(g),h)=c(a,k)$ is a relation of the
chunk $K_{A,E}$, so on a good fiber $f(g)f(h)$ and $f(k)$ differ on at most a fraction $\frac1{2r}$ of the points, and the bad fibers form a fraction at
most $\frac2{4r}$. So $f$ is an $r^{-1}$-morphism. For distinct $g,h\in E$: if $\pi(g)\ne\pi(h)$, then on every fiber $a$ with $a\pi(g),a\pi(h)\in A$ the
two images have different quotient coordinates, a fraction at least $1-\frac1{2r}$ of the points; if $\pi(g)=\pi(h)$, then $c(a,g)\ne c(a,h)$ are distinct
elements of $K_{A,E}$, so on the fibers with $a\pi(g)\in A$ the images differ on a fraction at least $1-\frac1{2r}$ of each fiber, and in total on at least
$(1-\frac1{4r})(1-\frac1{2r})\ge1-\frac1r$. So $f$ is $(1-r^{-1})$-expansive, on $|A|\,m$ points.
\end{proof}

When the cocycle takes finitely many values on $Q\times E$, for instance for direct products, $K_{A,E}$ lies in a fixed finite set and a bound with a
fixed kernel chunk follows, as Theorem~3.15 asserts. For $H_3$ with boxes $A=[0,M)^2$ and $M$ of order $r$, $K_{A,E_H}$ lies in the finite group of
permutations of the first $2M+O(1)$ points of ray~1, whose regular representation gives $\sigma_{K_{A,E_H}}\le(2M+O(1))!$; so the proposition gives
$\log\sigma_{E_H}(r)=O(r\log r)$, consistent with, and weaker than, Proposition~\ref{prop:sofic-upper}.

All numbering refers to the published version~\cite{Cornulier2013}. Theorem~3.15 and Example~3.16 read the same there
as in the arXiv version, where Proposition~3.17 appears as Example~3.17.

%======================================================================

%======================================================================
\section{Open problems}\label{sec:open}
%======================================================================

\begin{enumerate}[label=(\arabic*)]
\item \emph{The exponents.} The sofic exponent of $E_H$ in $r$ lies between $1/(\log130+1)\approx1/8.02$ and $1/2$ (up to the logarithm). For unitary
approximations with relator defect $e$ in the normalized Hilbert--Schmidt norm, $\log D$ lies between order $e^{-2\beta_H}$, with
$2\beta_H\approx0.133$ (Theorem~\ref{thm:unitary}), and $O(e^{-2/3})$ (Corollary~\ref{cor:unitary-upper}). The exponent is essentially
intrinsic to the particular substitution certificates used here: reweighting these fillings cannot go below $\log\rho\approx7.018$, with $\rho\approx129.64$ the spectral radius of their
substitution matrix, and $\log130$ is the certified choice. Since the area is at least quadratic
(Proposition~\ref{prop:area-lower}), the present uniform-area transport and sector argument cannot give more than $1/3$ in $r$ for permutations and $2/5$ in $1/e$
for unitaries. A quadratic Dehn function for $H_3$ would reach these values: every word in the argument has linear length. Theorem~\ref{thm:dehn}
gives the polynomial upper bound $O(n^{\log130+4})$, while for Thompson's group $F$ a quadratic Dehn function is a
theorem~\cite{Guba2006}. Shared fillings, improved aggregate defect estimates, or a stronger dimension argument are possible routes beyond these limits.
\item \emph{Permutation versus unitary approximation.} Relator errors add up in measure for permutations and in norm for unitaries, so the
two profiles may differ. The constructions here already have different exponents: in the normalized Hilbert--Schmidt norm, permutation models of $H_3$ give
$\log D=O(e^{-1})$ and unitary models $O(e^{-2/3})$ (Section~\ref{sec:unitary-models}). Can the unitary profile of a chunk of an amenable group be
polynomially smaller, in exponent, than its sofic profile?
\item \emph{Other groups.} Every Houghton group $H_m$ with $m\ge3$ contains $H_3$; the same method applies to groups with a displaced
$S_3$ structure and polynomially bounded far-commutator area. The proof of~\cite[Proposition~3.17]{Cornulier2013} for Baumslag--Solitar groups
uses the same extension theorem; whether its conclusion holds is open.
\item \emph{Extensions.} Proposition~\ref{prop:extension} bounds the profile of an extension by the Følner size of the quotient and the profile of a
kernel chunk that grows with $r$ through the cocycle of a section; the kernel chunk is fixed when the cocycle takes finitely many values. Which
extensions of groups with polynomial profile by amenable quotients have polynomial profile?
\end{enumerate}

\appendix
%======================================================================
\section{Fillings for the far-commutator bound}\label{app:area}
%======================================================================

Areas are taken in $P$; free conjugation and inversion do not change them, and freely $[x,yz]=[x,y]\cdot y[x,z]y^{-1}$.
Write $\tau_j=b^j\alpha b^{-j}$, $t=b^{-1}a$, $\rho=a^{-1}b^{-1}ab$, $t_i=a^ita^{-i}$, $X_n=b^{-n}a^n$, and
\[
A(d)=\mathrm{Area}[\sigma_0,\sigma_d],\quad E(i)=\mathrm{Area}[\sigma_i,t],\quad U(j)=\mathrm{Area}(\tau_j\sigma_j^{-1}),\quad
B(k)=\mathrm{Area}[\sigma_{-k},b].
\]
Since $b^{-n}(\tau_n\sigma_n^{-1})b^n=[\alpha,X_n]$ freely, $U(n)=\mathrm{Area}[\alpha,X_n]$. Since
$[x^{-1},y]=x^{-1}[y,x]x$, the signs in the definition of $A^*$ do not matter and $A^*(M)=\max_{2\le d\le M}A(d)$. Every
relation below with a numerical area is certified in the ancillary files (Section~\ref{sec:area}).

\paragraph{Base relations.} $[\alpha,t]$ is a rotation of $r_5^{-1}$; $aba^{-1}=\alpha b$ and $bab^{-1}=\alpha^{-1}a$ are
rotations of $r_4^{\pm1}$; $[\sigma_i,\sigma_{i+2}]$ is a conjugate of $r_3$; all have area $1$. The relation
$B_0:=\mathrm{Area}[\alpha,\sigma_1\sigma_0b]\le5$ follows by replacing $b\alpha^{-1}b^{-1}$ with $\sigma_1^{-1}$ (one $r_5$), which
leaves $\sigma_0\sigma_1\sigma_0\sigma_1^{-1}\sigma_0^{-1}\sigma_1^{-1}$, then three $r_1$ cells and one $r_2$ cell.

\begin{lemma}[Transfer and halving]\label{lem:transfer}
For $i\ge1$, $j\ge1$, $k\ge2$:
\begin{enumerate}[label=(\roman*)]
\item $E(i)\le2i+1+\sum_{d=2}^{i+1}A(d)$, and $E(0)=1$;
\item $U(j)\le\sum_{i=0}^{j-1}E(i)$;
\item $B(k)\le2k+5+\sum_{d=2}^{k-1}A(d)$;
\item $A(k+j)\le A(k)+2jB(k)+2U(j)$.
\end{enumerate}
\end{lemma}

\begin{proof}
(i) By $bab^{-1}=\alpha^{-1}a$, $t_k=t_{k+1}\sigma_{k-1}$ with one cell, so $t=t_i\sigma_{i-2}\cdots\sigma_0\sigma_{-1}$ with $i$ cells, and
$i$ more for $t^{-1}$. Then $[\sigma_i,t_iS]=[\sigma_i,t_i]\cdot t_i[\sigma_i,S]t_i^{-1}$ with $S=\sigma_{i-2}\cdots\sigma_{-1}$; the first
factor is $a^i[\alpha,t]a^{-i}$, and the second is a product of conjugates of $[\sigma_i,\sigma_l]$ for $l=i-2,\dots,-1$, of widths
$2,\dots,i+1$. (ii) Freely $\tau_{i+1}=b\tau_ib^{-1}$; replace $\tau_i$ by $\sigma_i$ at cost $U(i)$, and use that
$b\sigma_ib^{-1}\sigma_{i+1}^{-1}=b[\sigma_i,t]b^{-1}$ freely. (iii) Freely $[\sigma_{-k},b]=a^{-k}[\alpha,a^kba^{-k}]a^k$, and
$a^kba^{-k}=\sigma_{k-1}\cdots\sigma_1\sigma_0b$ with $k$ cells (twice, for the element and its inverse); then move $\alpha$ past
$\sigma_{k-1},\dots,\sigma_2$ and finish with $B_0$. (iv) Conjugating by $a^{-k}$ turns $[\sigma_0,\sigma_{k+j}]$ into
$[\sigma_{-k},\sigma_j]$. Replace $\sigma_j^{\pm1}$ by $\tau_j^{\pm1}$ ($2U(j)$), move $\sigma_{-k}^{\pm1}$ across the $2j$ letters $b^{\pm1}$
($2jB(k)$), and what remains is $b^j[\sigma_{-k},\alpha]b^{-j}$, a conjugate of $[\sigma_0,\sigma_k]$.
\end{proof}

The lemma gives, in this order: $A(2)=1$, $E(0)=U(1)=1$, $E(1)\le4$, $U(2)\le5$, $B(2)\le9$, $A(3)\le21$ ($k=2,j=1$),
$A(4)\le47$ ($k=j=2$), $E(2)\le27$, $U(3)\le32$, $B(3)\le12$, $A(5)\le79$ ($k=3,j=2$), $A(6)\le157$ ($k=j=3$) and $E(3)\le76$.
Part (iv) with $k=2$ gives $A(m)\le2U(m-2)+18m-35$ for $m\ge3$.

\paragraph{The cheap letters.} We use $g^h=h^{-1}gh$. Freely, $\psi(t)=b^{-2}a^2=t\rho t$, $t^b=b^{-2}ab=t\rho$, and
$\psi(\rho)=\rho^a\rho^{ba}\rho\rho^b=t^{-1}\rho^b\rho^{b^2}t\cdot\rho\cdot\rho^b$ (from $[x^2,y^2]=(xcx^{-1})(xy\,c\,y^{-1}x^{-1})\,c\,(ycy^{-1})$ with
$c=[x,y]$, $x=a^{-1}$, $y=b^{-1}$, and $a=bt$). The one relation needed is
\[
\rho^b=R':=t\rho t^{-1}\cdot\rho\cdot t\rho t^{-1},\qquad\text{area}\le14 .
\]
Replace each $\rho$ by $\alpha^{ba}$ (four $r_4$ cells); freely the relation becomes $\alpha^{bab}=\alpha^{b^2}\alpha^{ba}\alpha^{b^2}$, and after
conjugating by $(bab)^{-1}$ it reads $\alpha=\alpha^g\alpha^{g'}\alpha^g$ with $g=ba^{-1}b^{-1}$ and $g'=bab^{-1}a^{-1}b^{-1}$. Six cells of
$bab^{-1}=\alpha^{-1}a$ turn this into $\alpha=\alpha^{-1}\sigma_1\alpha\cdot b\alpha b^{-1}\cdot\alpha^{-1}\sigma_1\alpha$; one $r_5$ cell, two $r_1$ cells
and one $r_2$ cell finish. Also $[\alpha,t]$ has area $1$ and $[\alpha,\rho]$ area at most $5$: $\rho\alpha^{-1}\rho^{-1}$ becomes
$a^{-1}b^{-1}\sigma_2^{-1}ba$ after one $r_5$ cell, and $b^{-1}\sigma_2^{-1}b\to\sigma_1^{-1}$ costs $E(1)\le4$. Consequently, for a word
$W$ in $t^{\pm1},\rho^{\pm1}$, $W^b=W(t\rho,R')$ at area $14\,\#_\rho(W)$, and $\psi(W)=W(t\rho t,R)$ at area $84\,\#_\rho(W)$, where
\[
R=t^{-1}\cdot R'\cdot R'(t\rho,R')\cdot t\cdot\rho\cdot R'
\]
uses the relation six times per letter $\rho$; $|R'|=7$ and $|R|=46$.

\begin{lemma}[The doubling lifts]\label{lem:psi}
$\psi(\alpha)=\Pi:=\sigma_1\sigma_2\sigma_0\sigma_1$ with $6$ cells, and $\mathrm{Area}\,\psi(r_i)\le17,359,864,0,182$ for $i=1,\dots,5$.
Hence $\mathrm{Area}\,\psi(w)\le864\,\mathrm{Area}(w)$ for every null word $w$.
\end{lemma}

\begin{proof}
The identity for $[x^2,y^2]$ with $x=a$, $y=b$ and four $r_4$ cells ($[a,b]=\alpha$) give
$\psi(\alpha)=\sigma_1\cdot a(b\alpha b^{-1})a^{-1}\cdot\alpha\cdot b\alpha b^{-1}$, and two $r_5$ cells give $\Pi$. $\psi(r_4)$ is freely trivial.
$\psi(r_1)=\psi(\alpha)^2$: after two conversions, $\Pi^2$ costs five cells ($\sigma_1^2$, one $r_3$ swap, then $\sigma_2^2,\sigma_0^2,\sigma_1^2$).
$\psi(r_5)=a^2\psi(\alpha)a^{-2}\cdot b^2\psi(\alpha)^{-1}b^{-2}$: after two conversions, $a^2\Pi a^{-2}=\sigma_3\sigma_4\sigma_2\sigma_3$ freely, and
$b^2\sigma_i^{-1}b^{-2}\to\sigma_{i+2}^{-1}$ costs $E(i)+E(i+1)$; over the letters of $\Pi$ this is $E(0)+3E(1)+3E(2)+E(3)\le170$.
$\psi(r_3)=[\psi(\alpha),a^4\psi(\alpha)a^{-4}]$: after four conversions it is $[\Pi,\sigma_5\sigma_6\sigma_4\sigma_5]$, which splits freely into
sixteen conjugates of far commutators of widths $2,3,3,3,3,4,4,4,4,4,4,5,5,5,5,6$, costing $840$. $\psi(r_2)=(\psi(\alpha)\,a^2\psi(\alpha)a^{-2})^3$:
after six conversions it is the positive word $(\Pi\,\sigma_3\sigma_4\sigma_2\sigma_3)^3$ of length $24$ in $\sigma_0,\dots,\sigma_4$, which
represents the identity of $\Sym(6)$. Reduce it by the coset procedure: keep the word in the form $h\cdot c_k$ with
$c_k=\sigma_4\sigma_3\cdots\sigma_k$ and $h$ a word in $\sigma_0,\dots,\sigma_3$; absorbing a letter $\sigma_i$ costs far commutations for
$i\le k-2$, nothing for $i=k-1$, one cancellation for $i=k$, and far commutations plus one braid move for $i\ge k+1$, and then
recurse on $h$. A braid move $\sigma_i\sigma_{i-1}\sigma_i=\sigma_{i-1}\sigma_i\sigma_{i-1}$ costs four cells (three $r_1$, one $r_2$). Charging each
far commutation its actual width, the procedure uses $47$ moves at total cost $323$. The last claim follows by applying
$\psi$ to a product of conjugates of relators.
\end{proof}

\begin{proof}[Proof of Theorem~\ref{thm:area}]
Define words in $t^{\pm1},\rho^{\pm1}$ by $W_1=t$, $W_{2n}=W_n(t\rho t,R)$ and $W_{2n+1}=W_{2n}(t\rho,R')\,t$, and let
$\mathrm{Rw}(n)=\mathrm{Area}(X_nW_n^{-1})$. Then $\mathrm{Rw}(1)=0$. Since $X_{2n}=\psi(X_n)$ literally, Lemma~\ref{lem:psi} and the
cheap-letter relations give $\mathrm{Rw}(2n)\le864\,\mathrm{Rw}(n)+84\,|W_n|$ and $|W_{2n}|\le46|W_n|$. Since
$X_{2n+1}=X_{2n}^b\,t$ freely, $\mathrm{Rw}(2n+1)\le\mathrm{Rw}(2n)+14\cdot46|W_n|$ and $|W_{2n+1}|\le7|W_{2n}|+1$. Per binary digit,
$\mathrm{Rw}\le864\,\mathrm{Rw}'+728|W'|$ and $|W|\le323|W'|$, so with $s=\lfloor\log n\rfloor$,
\[
|W_n|\le323^s,\qquad \mathrm{Rw}(n)\le728\sum_{i<s}864^{s-1-i}323^i\le\frac{728}{541}\,864^s\le1.346\,n^{\log864}.
\]
Expanding $[\alpha,W_n]$ letter by letter, $U(n)\le2\,\mathrm{Rw}(n)+5|W_n|\le2.70\,n^{\log864}+5\,n^{\log323}$. For $m\ge5$,
$A(m)\le2U(m-2)+18m-35\le5.39\,m^{9.755}+10\,m^{8.336}+18m\le0.87\,m^{11}$, the ratio being largest at $m=5$; and
$A(2),A(3),A(4)\le1,21,47$.
\end{proof}

\begin{proof}[Proof of Theorem~\ref{thm:area-weighted}]
The proof above bounds every lifted cell by the largest area $864$ of the images $\psi(r_i)$. Count instead the cells of each type. Let
$C_{ij}$ be the number of cells of type $r_j$ in the certified filling of $\psi(r_i)$ from Lemma~\ref{lem:psi}; the certificates give
\[
C=\begin{pmatrix}4&0&1&8&4\\108&32&31&128&60\\204&68&64&384&144\\0&0&0&0&0\\36&12&14&84&36\end{pmatrix},
\]
with row sums $17,359,864,0,182$. For $w=(1,\frac{63}2,\frac{338}5,0,\frac{29}2)$, exact arithmetic gives $Cw\le130w$ and $C\mathbf1\le17w$
entrywise. The weight of $r_4$ is zero, but $\psi(r_4)$ is freely trivial, so cells of type $r_4$ disappear under $\psi$; the second
inequality pays for the ordinary cells created by the last application of $\psi$. Thus if a diagram has weighted count $S$ (a cell of
type $r_j$ counted with weight $w_j$), substituting the certified fillings gives weighted count at most $130S$ and ordinary area at most
$17S$. The relations $\rho^b=R'$ and $\psi(\rho)=R$ have ordinary costs $14$ and $84$ and weighted costs $48$ and $288$.

Let $T_n$ and $R_n$ count the letters $t^{\pm1}$ and $\rho^{\pm1}$ of $W_n$. The even step multiplies the column vector $(T,R)$ by
$D_0=\left(\begin{smallmatrix}2&26\\1&20\end{smallmatrix}\right)$, since $t\mapsto t\rho t$ and $R$ has $26$ letters $t^{\pm1}$ and $20$ letters
$\rho^{\pm1}$, and conjugation by $b$ multiplies it by $B_0=\left(\begin{smallmatrix}1&4\\1&3\end{smallmatrix}\right)$. Since
$(1,18)B_0D_0=(96,1654)\le96\,(1,18)$ and $(1,18)D_0\le96\,(1,18)$, the quantity $L_n=T_n+18R_n$ satisfies $L_n\le96L_{\lfloor n/2\rfloor}+1$,
so $L_n\le\frac{96}{95}96^s$ with $s=\lfloor\log n\rfloor$.

Let $S_n$ be the weighted count of the recursive filling of $X_nW_n^{-1}$ and $m=\lfloor n/2\rfloor$. Both branches give
$S_n\le130S_m+48T_m+1248R_m\le130S_m+\frac{208}3L_m$, hence $S_n\le\frac{3328}{1615}(130^s-96^s)\le\frac{21}{10}130^s$. For the ordinary count,
the second inequality applied to the diagram of level $m$ and the ordinary cost $14T_m+364R_m$ of the new relations give
$\mathrm{Rw}(n)\le17S_m+\frac{182}9L_m\le\frac9{20}130^s$, because $17\cdot\frac{21}{10}+\frac{182}9\cdot\frac{96}{95}<\frac9{20}\cdot130$. Expanding
$[\alpha,W_n]$ letter by letter, $U(n)\le2\,\mathrm{Rw}(n)+T_n+5R_n\le(\frac9{10}+\frac{96}{19})\,n^{\log130}$. Finally
$A(m)\le2U(m-2)+18m-35$ for $m\ge3$ gives $A^*(M)\le[2(\frac9{10}+\frac{96}{19})+18]M^{\log130}=\frac{2841}{95}M^{\log130}<30M^{\log130}$, and
$A(2)=1$.
\end{proof}

\section{From far commutators to the full Dehn function}\label{app:dehn}

\begin{proof}[Proof of Theorem~\ref{thm:dehn}]
Fix an integer $n\ge1$ and put $F_n=1+A^*(4n+2)$ and $U(0)=0$. Lemma~\ref{lem:transfer} gives
\[
U(q)\le\sum_{i=0}^{q-1}E(i)=O(n^2F_n)\qquad(0\le q\le2n).
\]
The implied constants in this proof are absolute.

\emph{Buffered collection.} Let $w$ be a null word of length at most $n$ and put $h=a^nb^n$.
The two quotient coordinates while reading $w$ after $h$ stay in $[0,2n]^2$. Write the current word as
$v\,s(m,q)$, where $s(m,q)=a^mb^q$ and $v$ is a word in the $\sigma_j^{\pm1}$.
Using $bab^{-1}=\alpha^{-1}a$ exactly $q$ times gives
$b^qab^{-q}=\tau_{q-1}^{-1}\cdots\tau_0^{-1}a$.
Replace each $\tau_j$ by $\sigma_j$ at cost $U(j)$. Thus
\[
\begin{split}
s(m,q)a&=(\sigma_{m+q-1}^{-1}\cdots\sigma_m^{-1})s(m+1,q),\\
s(m,q)a^{-1}&=(\sigma_{m-1}\cdots\sigma_{m+q-2})s(m-1,q)
\end{split}
\]
each costs at most $q+\sum_{j<q}U(j)=O(n^3F_n)$; an empty product is omitted.
The second rule is the inverse rearrangement of the first at $(m-1,q)$, and is used only when the next coordinate is nonnegative.
The rules $s(m,q)\alpha^{\pm1}=\sigma_{m+q}^{\pm1}s(m,q)$ cost $U(q)$, while multiplication by $b^{\pm1}$ merely changes $q$ and costs no cells.
The buffer ensures that all indices produced are nonnegative.

After at most $n$ rules the quotient coordinates return to $(n,n)$, and hence $hwh^{-1}$ has been changed to a word $v$
at cost $O(n^4F_n)$. It has at most $2n^2+n$ letters, with $0\le j\le4n$.
Since $w=1$, the word $v$ is the identity permutation on the first $4n+2$ points of ray~1. It remains to fill this permutation word
with a cost that is uniform in the size of that symmetric group.

\emph{Coxeter insertion.} In $S_{d+1}$, with adjacent transpositions $s_0,\ldots,s_{d-1}$, use the recursive normal form $u c_k$, where
$u\in S_d$, $c_k=s_{d-1}s_{d-2}\cdots s_k$ for $0\le k<d$, and $c_d=1$.
With right actions $c_k$ sends the last point to $k+1$; this determines $k$, after which $u$ is determined recursively.
To insert a generator on the right use
\[
c_ks_i=
\begin{cases}
s_ic_k,&i\le k-2,\\
c_{k-1},&i=k-1,\\
c_{k+1},&i=k,\\
s_{i-1}c_k,&i\ge k+1.
\end{cases}
\]
The first case uses far commutations, the middle cases concatenate or cancel, and the last case moves the final $s_i$ left across
$s_k,\ldots,s_{i-2}$, applies one braid, and moves the resulting $s_{i-1}$ left across $s_{i+1},\ldots,s_{d-1}$.
There are at most $d+1$ Coxeter moves at a level and at most $d$ levels per insertion, since only one generator is passed into $u$.
Consequently a word of length $l$ reduces to its unique normal form in $O(ld^2)$ moves, and the identity reduces to the empty word.

Apply this with $d=4n+1$ to $v$. First replace inverse letters by positive ones using $\sigma_j^2=1$, at one cell per letter.
A Coxeter square costs one cell in $P$, a braid costs at most four, and a far commutation costs at most $A^*(4n+2)$.
Thus the $O(n^2)$ letters of $v$ fill at cost $O(n^4F_n)$. Combining with collection and using invariance of area under free conjugation
proves the asserted bound for $\delta_P(n)$. Substituting Theorem~\ref{thm:area-weighted}, or Theorem~\ref{thm:area}, gives the two polynomial exponents.
\end{proof}

\section*{Use of AI assistance}
\addcontentsline{toc}{section}{Use of AI assistance}

Generative AI tools assisted with developing arguments and constructions, checking proofs and numerical identities, searching the literature,
writing verification code, and drafting and editing the manuscript. Automated review is not a substitute for mathematical proof. The accompanying
verification notes distinguish executable checks and kernel-checked Lean results from proofs in the manuscript that are not formalized in Lean.
The authors take responsibility for the content.

\begin{thebibliography}{99}
\bibitem{LiebRuskai1973} E.~H.~Lieb and M.~B.~Ruskai, Proof of the strong subadditivity of quantum-mechanical entropy, J. Math. Phys. \textbf{14} (1973), 1938--1941.
\bibitem{Cornulier2013} Y.~Cornulier, Sofic profile and computability of Cremona groups, Michigan Math. J. \textbf{62} (2013), 823--841; arXiv:1305.0993.
\bibitem{ElekSzabo2006} G.~Elek and E.~Szabó, On sofic groups, J. Group Theory \textbf{9} (2006), 161--171; arXiv:math/0305352.
\bibitem{Lee2012} S.~R.~Lee, Geometry of Houghton's groups, PhD thesis, University of Oklahoma (2012); arXiv:1212.0257.
\bibitem{Houghton1978} C.~H.~Houghton, The first cohomology of a group with permutation module coefficients, Arch. Math. \textbf{31} (1978), 254--258.
\bibitem{BCMR2014} J.~Burillo, S.~Cleary, A.~Martino and C.~E.~Röver, Commensurations and metric properties of Houghton's groups, Pacific J. Math. \textbf{285} (2016), 289--301; arXiv:1403.0026.
\bibitem{Slofstra2018} W.~Slofstra, A group with at least subexponential hyperlinear profile, arXiv:1806.05267 (2018).
\bibitem{SlofstraVidick2018} W.~Slofstra and T.~Vidick, Entanglement in non-local games and the hyperlinear profile of groups, Ann. Henri Poincaré \textbf{19} (2018), 2979--3005; arXiv:1711.10676.
\bibitem{ArzhantsevaCherix2020} G.~Arzhantseva and P.-A.~Cherix, Quantifying metric approximations of discrete groups, Ann. Fac. Sci. Toulouse Math. \textbf{33} (2024), 831--877; arXiv:2008.12954.
\bibitem{BCJM2023} P.~Burton, M.~Chaudkhari, K.~Juschenko and K.~Muliarchyk, Hyperlinear approximations to amenable groups come from sofic approximations, arXiv:2311.09202 (2023).
\bibitem{GowersHatami2017} W.~T.~Gowers and O.~Hatami, Inverse and stability theorems for approximate representations of finite groups, Sb. Math. \textbf{208} (2017), 1784--1817; arXiv:1510.04085.
\bibitem{DOT2019} M.~De Chiffre, N.~Ozawa and A.~Thom, Operator algebraic approach to inverse and stability theorems for amenable groups, Mathematika \textbf{65} (2019), 98--118; arXiv:1706.04544.
\bibitem{Guba2006} V.~S.~Guba, The Dehn function of Richard Thompson's group $F$ is quadratic, Invent. Math. \textbf{163} (2006), 313--342.
\bibitem{DouglasMghirbiB} S.~Douglas and N.~Mghirbi, Causal quantum-channel simulation: memory beyond entropy, companion manuscript (2026); \url{https://github.com/Apsiape/causal-quantum-memory}.
\end{thebibliography}

\end{document}
